\documentclass[dvipdfmx, xcolor=svgnames]{beamer}
%\documentclass[xcolor=svgnames]{beamer}

\usetheme{Madrid}
\setbeamertemplate{navigation symbols}{}
\setbeamertemplate{footline}[frame number]
\setbeamertemplate{footline}[section]
\usepackage{graphicx}
\usepackage{amsmath}
\usepackage{amssymb}
\usepackage{txfonts}
\usepackage{tabularx}
\usepackage{color}
%\mathversion{bold}
\renewcommand{\familydefault}{\sfdefault}
\renewcommand{\kanjifamilydefault}{\gtdefault}
\setbeamerfont{title}{size=\large,series=\bfseries}
\setbeamerfont{frametitle}{size=\large,series=\bfseries}
%\setbeamerfont{framesubtitle}{size=\small,series=\bfseries}
\setbeamercolor{title}{fg=white, bg=Navy}
\setbeamertemplate{frametitle}[default][left]
\setbeamercolor{frametitle}{fg=white, bg=Navy}
%\setbeamercolor{framesubtitle}{fg=white, bg=black} %%%%%%%%%% 反映されないよ
\usefonttheme{professionalfonts}

\usepackage[deluxe]{otf}
\usepackage[noalphabet]{pxchfon}
\setboldgothicfont{HaranoAjiGothic-Medium.otf}

\theoremstyle{plain}
%\setbeamertemplate{theorems}
%\newtheorem{main}{Main Result}
\newtheorem{thm}{Theorem}[section]
\newtheorem{cor}{Corollary}[section]
\newtheorem{prop}{Proposition}[section]
\newtheorem{lemm}{Lemma}[section]
\theoremstyle{definition}
\newtheorem*{defi}{Definition}
%\newtheorem{fact}{Fact}[section]
\newtheorem{question}{Question}
\newtheorem{exam}{Example}[section]
\newtheorem{assumption}{Assumption}[section]
%\numberwithin{equation}{section}
\theoremstyle{remark}
\newtheorem{remark}{Remark}[section]


\usepackage{amsmath,amssymb,amsthm,amscd,ascmac,mathrsfs,bm,type1cm,multirow,array,enumerate,mathrsfs,tabularx}
\usepackage{wrapfig,float}
% \usepackage{graphicx}
\usepackage{color}
\usepackage{arydshln} % 行列 c:c 
% \usepackage[dvipdfmx]{graphicx}

\renewcommand{\arraystretch}{1.1}

\newcounter{main} % カウンタの宣言
\setcounter{main}{0} % カウンタの初期化
\resetcounteronoverlays{main}
\newcommand{\useMain}[1][]{\refstepcounter{main}{#1}Main Result \,{\themain}}


\newcounter{Prop} % カウンタの宣言
\setcounter{Prop}{0} % カウンタの初期化
\resetcounteronoverlays{Prop}
\newcommand{\useProp}[1][]{\refstepcounter{Prop}{#1}Proposition \,{\theProp}}



\def\rnum#1{\expandafter{\romannumeral #1}} 
\def\Rnum#1{\uppercase\expandafter{\romannumeral #1}} 
\def\hsymb#1{\mbox{\strut\rlap{\smash{\Huge$#1$}}\quad}} % 行列の省略記号*

\newcommand{\BigFig}[1]{\parbox{12pt}{\Huge #1}}
\newcommand{\BigZero}{\BigFig{0}} 


\setcounter{tocdepth}{1}


\title{研究内容の紹介（20分）}
\author{東京理科大学理学研究科 \quad
        小林 穂乃香}
%\institute{小池研究室}
\date{山口大学面接\, 於\\
      2022年9月5日}

\begin{document}
%%%%% page 1 (title page) %%%%%
\frame{\titlepage}




\section{擬リーマン多様体とは}

\frame{
\frametitle{\insertsection}
\hspace{10pt}$g$\hspace{10pt} : \(M\)上の対称$(0,2)$テンソル場\\
\hspace{5pt}$X,\, Y\, \in TM$ \vspace{10pt}\\
\underline{$(M, g)$ : リーマン多様体}\vspace{-5pt}
\begin{table}
\begin{tabular}{lll}
\hspace{-3pt}  
$g$ : $M$上のリーマン計量 
\hspace{10pt}
&
\uncover<1->{\hspace{4pt} i. e.}
&
\uncover<1->{$g(X,X)\geq0$ {\it \, for all $X$, and}}
\\
 \hspace{48pt}\(\cdots \)正定値
&
&
\uncover<1->{$g(X,X)=0 \Leftrightarrow X=0$}
\end{tabular}
\vspace{6pt}
\end{table}
\uncover<1->{
\underline{$(M, g)$ : 擬リーマン多様体}\vspace{-5pt}
\begin{table}
  \hspace{20pt}
\begin{tabular}{lll}
  \hspace{-24pt}
$g$ : $M$上の擬リーマン計量 
&
\uncover<1->{\hspace{8pt}i. e. }
&
\uncover<1->{
  \hspace{-9pt}
$g(X,Y)=0$\, {\it for all}\, $Y\in TM$ 
}
\\
\uncover<1->{\hspace{28pt}\(\cdots \) 非退化}
&
&
\uncover<1->{\hspace{30pt} $\Rightarrow X=0$}
\end{tabular}
\end{table}

\uncover<2->{
  \begin{defi}{}
    \begin{table}
    \begin{center}
    \begin{tabular}{lllcl}
    $X$ : \textcolor{red}{{\it spacelike}} & 
    $\overset{\mathrm{def}}{\Leftrightarrow}$ &
    $\langle X, X\rangle >0$ &
    or &
    $X=0$
    \\
    $X$ : \textcolor{red}{{\it timelike}} & 
    $\overset{\mathrm{def}}{\Leftrightarrow}$ & 
    $\langle X, X\rangle <0$ & 
    \\
    $X$ : \textcolor{red}{{\it null}} & 
    $\overset{\mathrm{def}}{\Leftrightarrow}$ &
    $\langle X, X\rangle =0$ &
    and &
    $X\neq 0$
    \end{tabular}
    \end{center}
    \vspace{-4pt}
    \end{table}
    \end{defi}
}
}}






\frame{
  \frametitle{\insertsection}
${\bf x}=(x_1,\ldots ,x_{m+1})$, \, ${\bf y}=(y_1,\ldots,y_{m+1})$ \, $\in \mathbb{E}^{m+1}_{\textcolor<2->{red}{s}}$
\hspace{5pt} 
\uncover<2->{
  \begin{minipage}{0.3\textwidth}
    {\small 
    \(
      \left\{
    \begin{array}{ll}
    s=0 \, \text{: リーマン多様体}\\
    s=1 \, \text{: ローレンツ多様体}
    \end{array}
      \right.
    \)
    }
  \end{minipage}
}
\vspace{5pt}

\hspace{20pt}
\(\langle {\bf x},{\bf y}\rangle :=\textcolor<2->{red}{-}\sum_{i=1}^{\textcolor<2->{red}{s}}x_iy_i +\sum_{j=\textcolor<2->{red}{s+1}}^{m+1}x_jy_j\)
\hspace{20pt}
\uncover<2->{\textcolor{red}{$s$ : 指数}}
\vspace{0pt}
\begin{table}
\hspace{0pt}
\begin{tabular}{rcl}
\hspace{-15pt}
\begin{minipage}[l]{0.4\textwidth}
\uncover<3->{\textcolor{black}{\underline{リーマン多様体}}}
\hspace{-3pt}
\vspace{5pt}
\end{minipage}
&
\hspace{-3pt}
&
\hspace{-20pt}
\begin{minipage}[l]{0.4\textwidth}
%\textcolor{black}{
\uncover<3->{\underline{擬リーマン多様体}}
%}
\vspace{5pt}
\end{minipage}
\\
\hspace{-15pt}
\begin{minipage}[l]{0.4\textwidth}
\shortstack{
%\textcolor{black}{$\mathbb{S}^2_1$ : pseudo-sphere in $\mathbb{R}^3_1$}
\uncover<3->{${\small \mathbb{S}^2:=\{ {\bf x}\in\mathbb{E}^3\,} {\large \vert} {\small \, \langle {\bf x},{\bf x}\rangle =1 \} }$}\\
\hspace{5pt}\uncover<4->{${\small \mathbb{H}^2:=\{ {\bf x}\in\mathbb{E}^3_1\,} {\large \vert} {\small \, \langle {\bf x},{\bf x}\rangle =-1 \} }$}
}
\end{minipage}
\vspace{7pt}
&
&
\hspace{-20pt}
\begin{minipage}[l]{0.4\textwidth}
\shortstack{
%\textcolor{black}{$\mathbb{S}^2_1$ : pseudo-sphere in $\mathbb{R}^3_1$}
\uncover<3->{${\small \mathbb{S}^2_1:=\{ {\bf x}\in\mathbb{E}^3_1\,} {\large \vert} {\small \, \langle {\bf x},{\bf x}\rangle =1 \} }$}\\
\hspace{5pt}\uncover<4->{${\small \mathbb{H}^2_1:=\{ {\bf x}\in\mathbb{E}^3_2\,} {\large \vert} {\small \, \langle {\bf x},{\bf x}\rangle =-1 \} }$}
}
\end{minipage}
\\
\begin{minipage}[c]{0.4\textwidth}
\hspace{10pt}
\uncover<3->{
\scalebox{0.5}{%WinTpicVersion4.32a
{\unitlength 0.1in%
\begin{picture}(28.6000,26.5000)(4.2000,-30.0000)%
% STR 2 0 3 0 Black White  
% 4 1840 1777 1840 1790 4 2400 0 0
% O
\put(18.4000,-17.9000){\makebox(0,0)[rt]{O}}%
% VECTOR 2 0 3 0 Black White  
% 2 1850 3000 1850 350
% 
\special{pn 8}%
\special{pa 1850 3000}%
\special{pa 1850 350}%
\special{fp}%
\special{sh 1}%
\special{pa 1850 350}%
\special{pa 1830 417}%
\special{pa 1850 403}%
\special{pa 1870 417}%
\special{pa 1850 350}%
\special{fp}%
% VECTOR 2 0 3 0 Black White  
% 2 420 1780 3280 1780
% 
\special{pn 8}%
\special{pa 420 1780}%
\special{pa 3280 1780}%
\special{fp}%
\special{sh 1}%
\special{pa 3280 1780}%
\special{pa 3213 1760}%
\special{pa 3227 1780}%
\special{pa 3213 1800}%
\special{pa 3280 1780}%
\special{fp}%
% FUNC 2 0 3 0 Black White  
% 10 420 350 3280 3000 1850 1780 2950 1780 1850 680 420 350 3280 3000 50 2 0 3 1 0
% ///sin(t)^2+cos(t)^2///0///2pi
\special{pn 8}%
\special{pa 2950 1780}%
\special{pa 2950 1752}%
\special{pa 2949 1738}%
\special{pa 2949 1725}%
\special{pa 2946 1683}%
\special{pa 2944 1669}%
\special{pa 2943 1656}%
\special{pa 2939 1628}%
\special{pa 2937 1615}%
\special{pa 2933 1587}%
\special{pa 2930 1574}%
\special{pa 2928 1560}%
\special{pa 2925 1547}%
\special{pa 2922 1533}%
\special{pa 2919 1520}%
\special{pa 2915 1506}%
\special{pa 2912 1493}%
\special{pa 2908 1480}%
\special{pa 2904 1466}%
\special{pa 2892 1427}%
\special{pa 2887 1414}%
\special{pa 2883 1401}%
\special{pa 2868 1362}%
\special{pa 2862 1349}%
\special{pa 2857 1337}%
\special{pa 2851 1324}%
\special{pa 2845 1312}%
\special{pa 2839 1299}%
\special{pa 2833 1287}%
\special{pa 2827 1274}%
\special{pa 2820 1262}%
\special{pa 2814 1250}%
\special{pa 2779 1190}%
\special{pa 2771 1179}%
\special{pa 2764 1167}%
\special{pa 2748 1145}%
\special{pa 2740 1133}%
\special{pa 2732 1122}%
\special{pa 2723 1111}%
\special{pa 2715 1100}%
\special{pa 2706 1089}%
\special{pa 2697 1079}%
\special{pa 2689 1068}%
\special{pa 2680 1058}%
\special{pa 2670 1047}%
\special{pa 2652 1027}%
\special{pa 2642 1017}%
\special{pa 2633 1007}%
\special{pa 2623 997}%
\special{pa 2613 988}%
\special{pa 2603 978}%
\special{pa 2593 969}%
\special{pa 2583 959}%
\special{pa 2572 950}%
\special{pa 2562 941}%
\special{pa 2551 932}%
\special{pa 2540 924}%
\special{pa 2530 915}%
\special{pa 2519 907}%
\special{pa 2508 898}%
\special{pa 2496 890}%
\special{pa 2474 874}%
\special{pa 2462 866}%
\special{pa 2451 859}%
\special{pa 2439 851}%
\special{pa 2428 844}%
\special{pa 2368 809}%
\special{pa 2355 803}%
\special{pa 2331 791}%
\special{pa 2318 785}%
\special{pa 2306 779}%
\special{pa 2293 773}%
\special{pa 2280 768}%
\special{pa 2268 762}%
\special{pa 2229 747}%
\special{pa 2216 743}%
\special{pa 2203 738}%
\special{pa 2177 730}%
\special{pa 2163 726}%
\special{pa 2137 718}%
\special{pa 2123 715}%
\special{pa 2110 711}%
\special{pa 2097 708}%
\special{pa 2083 705}%
\special{pa 2070 702}%
\special{pa 2056 699}%
\special{pa 2042 697}%
\special{pa 2029 695}%
\special{pa 2015 692}%
\special{pa 2001 690}%
\special{pa 1988 689}%
\special{pa 1974 687}%
\special{pa 1960 686}%
\special{pa 1946 684}%
\special{pa 1933 683}%
\special{pa 1905 681}%
\special{pa 1891 681}%
\special{pa 1878 680}%
\special{pa 1822 680}%
\special{pa 1808 681}%
\special{pa 1795 681}%
\special{pa 1753 684}%
\special{pa 1739 686}%
\special{pa 1726 687}%
\special{pa 1698 691}%
\special{pa 1685 693}%
\special{pa 1657 697}%
\special{pa 1644 700}%
\special{pa 1630 702}%
\special{pa 1617 705}%
\special{pa 1603 708}%
\special{pa 1590 711}%
\special{pa 1576 715}%
\special{pa 1563 718}%
\special{pa 1550 722}%
\special{pa 1536 726}%
\special{pa 1497 738}%
\special{pa 1484 743}%
\special{pa 1471 747}%
\special{pa 1432 762}%
\special{pa 1419 768}%
\special{pa 1407 773}%
\special{pa 1394 779}%
\special{pa 1382 785}%
\special{pa 1369 791}%
\special{pa 1357 797}%
\special{pa 1344 803}%
\special{pa 1332 810}%
\special{pa 1320 816}%
\special{pa 1260 851}%
\special{pa 1249 859}%
\special{pa 1237 866}%
\special{pa 1215 882}%
\special{pa 1203 890}%
\special{pa 1192 898}%
\special{pa 1181 907}%
\special{pa 1170 915}%
\special{pa 1159 924}%
\special{pa 1149 933}%
\special{pa 1138 941}%
\special{pa 1128 950}%
\special{pa 1117 960}%
\special{pa 1097 978}%
\special{pa 1087 988}%
\special{pa 1077 997}%
\special{pa 1067 1007}%
\special{pa 1058 1017}%
\special{pa 1048 1027}%
\special{pa 1039 1037}%
\special{pa 1029 1047}%
\special{pa 1020 1058}%
\special{pa 1011 1068}%
\special{pa 1002 1079}%
\special{pa 994 1090}%
\special{pa 985 1100}%
\special{pa 977 1111}%
\special{pa 968 1122}%
\special{pa 960 1134}%
\special{pa 944 1156}%
\special{pa 936 1168}%
\special{pa 929 1179}%
\special{pa 921 1191}%
\special{pa 914 1202}%
\special{pa 879 1262}%
\special{pa 873 1275}%
\special{pa 861 1299}%
\special{pa 855 1312}%
\special{pa 849 1324}%
\special{pa 843 1337}%
\special{pa 838 1350}%
\special{pa 832 1362}%
\special{pa 817 1401}%
\special{pa 813 1414}%
\special{pa 808 1427}%
\special{pa 800 1453}%
\special{pa 796 1467}%
\special{pa 788 1493}%
\special{pa 785 1507}%
\special{pa 781 1520}%
\special{pa 778 1533}%
\special{pa 775 1547}%
\special{pa 772 1560}%
\special{pa 769 1574}%
\special{pa 767 1588}%
\special{pa 765 1601}%
\special{pa 762 1615}%
\special{pa 760 1629}%
\special{pa 759 1642}%
\special{pa 757 1656}%
\special{pa 756 1670}%
\special{pa 754 1684}%
\special{pa 753 1697}%
\special{pa 751 1725}%
\special{pa 751 1739}%
\special{pa 750 1752}%
\special{pa 750 1808}%
\special{pa 751 1822}%
\special{pa 751 1835}%
\special{pa 754 1877}%
\special{pa 756 1891}%
\special{pa 757 1904}%
\special{pa 761 1932}%
\special{pa 763 1945}%
\special{pa 767 1973}%
\special{pa 770 1986}%
\special{pa 772 2000}%
\special{pa 775 2013}%
\special{pa 778 2027}%
\special{pa 781 2040}%
\special{pa 785 2054}%
\special{pa 788 2067}%
\special{pa 792 2080}%
\special{pa 796 2094}%
\special{pa 808 2133}%
\special{pa 813 2146}%
\special{pa 817 2159}%
\special{pa 832 2198}%
\special{pa 838 2211}%
\special{pa 843 2223}%
\special{pa 849 2236}%
\special{pa 855 2248}%
\special{pa 861 2261}%
\special{pa 867 2273}%
\special{pa 873 2286}%
\special{pa 880 2298}%
\special{pa 886 2310}%
\special{pa 921 2370}%
\special{pa 929 2381}%
\special{pa 936 2393}%
\special{pa 952 2415}%
\special{pa 960 2427}%
\special{pa 968 2438}%
\special{pa 977 2449}%
\special{pa 985 2460}%
\special{pa 994 2471}%
\special{pa 1003 2481}%
\special{pa 1011 2492}%
\special{pa 1020 2502}%
\special{pa 1030 2513}%
\special{pa 1048 2533}%
\special{pa 1058 2543}%
\special{pa 1067 2553}%
\special{pa 1077 2563}%
\special{pa 1087 2572}%
\special{pa 1097 2582}%
\special{pa 1107 2591}%
\special{pa 1117 2601}%
\special{pa 1128 2610}%
\special{pa 1138 2619}%
\special{pa 1149 2628}%
\special{pa 1160 2636}%
\special{pa 1170 2645}%
\special{pa 1181 2653}%
\special{pa 1192 2662}%
\special{pa 1204 2670}%
\special{pa 1226 2686}%
\special{pa 1238 2694}%
\special{pa 1249 2701}%
\special{pa 1261 2709}%
\special{pa 1272 2716}%
\special{pa 1332 2751}%
\special{pa 1345 2757}%
\special{pa 1369 2769}%
\special{pa 1382 2775}%
\special{pa 1394 2781}%
\special{pa 1407 2787}%
\special{pa 1420 2792}%
\special{pa 1432 2798}%
\special{pa 1471 2813}%
\special{pa 1484 2817}%
\special{pa 1497 2822}%
\special{pa 1523 2830}%
\special{pa 1537 2834}%
\special{pa 1563 2842}%
\special{pa 1577 2845}%
\special{pa 1590 2849}%
\special{pa 1603 2852}%
\special{pa 1617 2855}%
\special{pa 1630 2858}%
\special{pa 1644 2861}%
\special{pa 1658 2863}%
\special{pa 1671 2865}%
\special{pa 1685 2868}%
\special{pa 1699 2870}%
\special{pa 1712 2871}%
\special{pa 1726 2873}%
\special{pa 1740 2874}%
\special{pa 1754 2876}%
\special{pa 1767 2877}%
\special{pa 1795 2879}%
\special{pa 1809 2879}%
\special{pa 1822 2880}%
\special{pa 1878 2880}%
\special{pa 1892 2879}%
\special{pa 1905 2879}%
\special{pa 1947 2876}%
\special{pa 1961 2874}%
\special{pa 1974 2873}%
\special{pa 2002 2869}%
\special{pa 2015 2867}%
\special{pa 2043 2863}%
\special{pa 2056 2860}%
\special{pa 2070 2858}%
\special{pa 2083 2855}%
\special{pa 2097 2852}%
\special{pa 2110 2849}%
\special{pa 2124 2845}%
\special{pa 2137 2842}%
\special{pa 2150 2838}%
\special{pa 2164 2834}%
\special{pa 2203 2822}%
\special{pa 2216 2817}%
\special{pa 2229 2813}%
\special{pa 2268 2798}%
\special{pa 2281 2792}%
\special{pa 2293 2787}%
\special{pa 2306 2781}%
\special{pa 2318 2775}%
\special{pa 2331 2769}%
\special{pa 2343 2763}%
\special{pa 2356 2757}%
\special{pa 2368 2750}%
\special{pa 2380 2744}%
\special{pa 2440 2709}%
\special{pa 2451 2701}%
\special{pa 2463 2694}%
\special{pa 2485 2678}%
\special{pa 2497 2670}%
\special{pa 2508 2662}%
\special{pa 2519 2653}%
\special{pa 2530 2645}%
\special{pa 2541 2636}%
\special{pa 2551 2627}%
\special{pa 2562 2619}%
\special{pa 2572 2610}%
\special{pa 2583 2600}%
\special{pa 2603 2582}%
\special{pa 2613 2572}%
\special{pa 2623 2563}%
\special{pa 2633 2553}%
\special{pa 2642 2543}%
\special{pa 2652 2533}%
\special{pa 2661 2523}%
\special{pa 2671 2513}%
\special{pa 2680 2502}%
\special{pa 2689 2492}%
\special{pa 2698 2481}%
\special{pa 2706 2470}%
\special{pa 2715 2460}%
\special{pa 2723 2449}%
\special{pa 2732 2438}%
\special{pa 2740 2426}%
\special{pa 2756 2404}%
\special{pa 2764 2392}%
\special{pa 2771 2381}%
\special{pa 2779 2369}%
\special{pa 2786 2358}%
\special{pa 2821 2298}%
\special{pa 2827 2285}%
\special{pa 2839 2261}%
\special{pa 2845 2248}%
\special{pa 2851 2236}%
\special{pa 2857 2223}%
\special{pa 2862 2210}%
\special{pa 2868 2198}%
\special{pa 2883 2159}%
\special{pa 2887 2146}%
\special{pa 2892 2133}%
\special{pa 2900 2107}%
\special{pa 2904 2093}%
\special{pa 2912 2067}%
\special{pa 2915 2053}%
\special{pa 2919 2040}%
\special{pa 2922 2027}%
\special{pa 2925 2013}%
\special{pa 2928 2000}%
\special{pa 2931 1986}%
\special{pa 2933 1972}%
\special{pa 2935 1959}%
\special{pa 2938 1945}%
\special{pa 2940 1931}%
\special{pa 2941 1918}%
\special{pa 2943 1904}%
\special{pa 2944 1890}%
\special{pa 2946 1876}%
\special{pa 2947 1863}%
\special{pa 2949 1835}%
\special{pa 2949 1821}%
\special{pa 2950 1808}%
\special{pa 2950 1780}%
\special{fp}%
% VECTOR 2 0 3 0 Black White  
% 2 530 2230 3180 1340
% 
\special{pn 8}%
\special{pa 530 2230}%
\special{pa 3180 1340}%
\special{fp}%
\special{sh 1}%
\special{pa 3180 1340}%
\special{pa 3110 1342}%
\special{pa 3129 1357}%
\special{pa 3123 1380}%
\special{pa 3180 1340}%
\special{fp}%
% ELLIPSE 2 0 3 0 Black White  
% 4 1850 1780 2940 1940 2940 1940 2940 1940
% 
\special{pn 8}%
\special{ar 1850 1780 1090 160 0.0000000 6.2831853}%
\end{picture}}%
}
}
\end{minipage}
&
\begin{minipage}[c]{0.15\textwidth}
\hspace{-12pt}
\uncover<4->{\textcolor{red}{$\mathbb{S}^n_s$と$\mathbb{H}^n_{n-s}$は反等長}}
\hspace{-3pt}
\end{minipage}
&
\begin{minipage}[c]{0.4\textwidth}
\hspace{-12pt}
\uncover<3->{
\scalebox{0.4}{\input{pic_pseudo-sphere-8}}
}
\end{minipage}
\end{tabular}
\end{table}
}





\frame{
\frametitle{\insertsection}
  \begin{table}
   \hspace{-10pt}
    \begin{tabular}{cl}
      1905 & {\bf 特殊相対性理論}（A. Einstein）\\
           & この世界を, 空間\(3\)次元と時間\(1\)次元の"時空"として考える\\
           & 正定値とは限らない計量を持つ多様体を導入\\
           & \\
      1908 & {\bf ミンコフスキー幾何学}（H. Minkowski）\\
           & 特殊相対性理論を幾何学として再構成\\
           & 時空は\(4\)次元の空間として記述される \\
           & \\
      1915-1916& {\bf 一般相対性理論}（A. Einstein）\\
           & 重力を時空の曲がりとして捉える, リーマン幾何学を応用        
    \end{tabular}
  \end{table}
  \begin{center}
    \raisebox{-3mm}[0mm]{\includegraphics[width=15pt, height=20pt, scale=0.2]{yazirusitate.png}} 
    \vspace{7pt}

    {\bf 擬リーマン幾何学}
    \end{center}
}



% \begin{frame}{\insertsection}
% \underline{例2．} \quad 重力場における光の軌跡
% \vspace{-5pt}
% \begin{center}
% \scalebox{0.25}{\includegraphics{logoRikadai.png}}   
% \end{center}

% {\bf \underline{aの位置にある天体が，b地点から見える（重力レンズ効果）}}

% \begin{itemize}
%       \item 光は常に“真っ直ぐ”進む
%       \item 光の粒子の軌跡はnull測地線
% \end{itemize}
% \end{frame}



%   \begin{frame}[noframenumbering]{\insertsection}
%     \underline{例2．} \quad 重力場における光の軌跡
%     \vspace{-5pt}
%     \begin{center}
%     \scalebox{0.13}{\includegraphics{logoRikadai3.jpg}}     
%     \end{center}
    
%     {\bf \underline{aの位置にある天体が，b地点から見える（重力レンズ効果）}}
    
%     \begin{itemize}
%           \item 光は常に“真っ直ぐ”進む
%           \item 光の粒子の軌跡はnull測地線
%     \end{itemize}
%     \end{frame}


\section{本日の発表の流れ}


\frame{
  \frametitle{\insertsection}
$(\, \tilde{M}, \langle \,,\rangle \,)$ : 擬球面 $\mathbb{S}^m_s$ \,または\, 擬双曲空間
$\mathbb{H}^m_s$　\vspace{5pt} 
\\
$(\, M, \langle \,, \rangle \,)$ : $\tilde{M}$の擬リーマン曲面,
\\
\hspace{45pt}
% \uncover<2->{
\textcolor<1-5>{red}{形作用素が対角化不可能, 平均曲率とスカラー曲率が一定}
  % }
\vspace{3pt} 
\\
\hspace{15pt} $\gamma$ \hspace{16pt} : $\tilde{M}$のnull曲線 \quad 
i.e. \(\langle \dot{\gamma}, \dot{\gamma}\rangle=0\)\, かつ\, \(\dot{\gamma} \neq 0\)
\vspace{20pt}
\pause
\\
{\large \underline{\, Part 1\, \(\cdots \) 擬双曲的ガウス写像による分類}}\vspace{2pt}
\begin{itemize}
\item \textcolor<6->{red}{B-scrollまたはcomplex circleの擬双曲的ガウス写像}
%\item 平行曲面の擬双曲的ガウス写像
% \item B-scrollの一般次元版であるgeneralized umbilical hypersurfaceの\\擬双曲的ガウス写像　 
\vspace{10pt}
\end{itemize}
\pause

{\large \underline{\, Part 2 \, \(\cdots \) generalizations of B-scroll in \(\tilde{M}^{m}_s\)}}\vspace{2pt}
\begin{itemize}
\item generalized umbilical hypersurface in \(\tilde{M}^{n+1}_1\) \, \(\cdots\) 次元一般化 
\pause
\item \textcolor<6->{red}{generalized umbilical hypersurface in \(\tilde{M}^{n+1}_2\) \, \(\cdots\) 指数及び次元一般化} 
\pause
\item \textcolor<6->{red}{generalized B-scroll\, in \, \(\tilde{M}^5_2\) \, \(\cdots\) 指数及び余次元一般化}
\end{itemize}
}




\section{Part 1-1 \quad 擬双曲的ガウス写像による分類}




% \frame{
%   \begin{center}
%     {\Large {\bf Part \Rnum{1}}}
%   \end{center}
% \begin{beamercolorbox}[
%   %wd=60mm, sep=2pt,
%   %shadow=true, 
%   rounded=true
%   ]{frametitle}
%   \begin{center}
%     {\Large 擬双曲的ガウス写像による分類}
%   \end{center}
%   \end{beamercolorbox}
% }



\subsection{研究背景}


\frame{
\frametitle{\insertsection \quad \quad{\normalsize $\cdots$ \insertsubsection}}
\({\bf x} : M \hookrightarrow \mathbb{S}^m\) : 等長はめ込み
\vspace{3pt}

\hspace{3pt}\((e_1^p, \ldots, e_n^p)\)\hspace{3pt} : \(M\)の向きと適合する\(T_pM\)の正規直交フレーム
\begin{table}
 \begin{tabular}{lll}
   ガウス写像\(\nu\)
 & \(\overset{\mathrm{def}}{\Leftrightarrow}\)
 & \(\nu(p) := e_1^p \wedge \cdots \wedge e_n^p\)\\
 球面的ガウス写像\(\tilde{\nu}\)
 & \(\overset{\mathrm{def}}{\Leftrightarrow}\)
 & \(\tilde{\nu}(p) := \textcolor{red}{{\bf x}(p)} \wedge e_1^p \wedge \cdots \wedge e_n^p \)
 \end{tabular}
 \end{table}
 \vspace{10pt}
 \pause

$\phi: M\rightarrow \mathbb{S}^{m}_s\subset \mathbb{E}^{m+1}_s$ (or $ \mathbb{H}^{m}_{s}\subset \mathbb{E}^{m+1}_{s+1}$) : $C^\infty$級写像\\
\begin{table}
\begin{center}
\begin{tabular}{lll}
$\phi$ : \(k\)-type & $\overset{\mathrm{def}}{\Leftrightarrow}$ & $\phi=\phi_1+ \phi_2 +\cdots +\phi_k$,\quad $\Delta \phi_i=\lambda_i\phi_i$ \, ($\lambda_i\in \mathbb{R}$)
\end{tabular}
\end{center}
\end{table}
\vspace{10pt}
\pause

\underline{部分多様体の平均曲率ベクトル場\(H\)と部分多様体のtype numberの関係}
\begin{table}
  \hspace{-10pt}
   \begin{tabular}{cl}
     1970年代 & ・\(\Delta H=\lambda H\)となるリーマン部分多様体のtype number\\
          & ・null \(2\)-typeかつ \(H\) : const\, なローレンツ曲面の完全分類\vspace{5pt}\\
     1980年代 & ・部分多様体のガウス写像のtype number \vspace{5pt}\\
     2007年 & ・球面にはめ込まれた部分多様体の球面的ガウス写像   
   \end{tabular}
   \vspace{10pt}
 \end{table}
}



% \frame{
%   \frametitle{\insertsection \quad \quad{\normalsize $\cdots$ \insertsubsection}}
%   \underline{擬双曲的ガウス写像}
% \begin{center}
% \scalebox{0.8}{%WinTpicVersion3.08
\unitlength 0.1in
\begin{picture}( 42.6000, 34.6000)(  5.5000,-37.6000)
% STR 2 0 3 0
% 3 3090 2223 3090 2240 4 3760
% O
\put(30.9000,-22.4000){\makebox(0,0)[rt]{O}}%
% VECTOR 3 0 3 0
% 2 3100 3760 3100 580
% 
\special{pn 4}%
\special{pa 3100 3760}%
\special{pa 3100 580}%
\special{fp}%
\special{sh 1}%
\special{pa 3100 580}%
\special{pa 3080 648}%
\special{pa 3100 634}%
\special{pa 3120 648}%
\special{pa 3100 580}%
\special{fp}%
% VECTOR 3 0 3 0
% 2 1380 2230 4780 2230
% 
\special{pn 4}%
\special{pa 1380 2230}%
\special{pa 4780 2230}%
\special{fp}%
\special{sh 1}%
\special{pa 4780 2230}%
\special{pa 4714 2210}%
\special{pa 4728 2230}%
\special{pa 4714 2250}%
\special{pa 4780 2230}%
\special{fp}%
% FUNC 2 0 3 0
% 9 1380 580 4780 3760 3100 2230 3700 2230 3100 1630 1380 580 4780 3760 50 3 0 2
% (exp(t)+exp(-t))/2///(exp(t)-exp(-t))/2///-10///10
\special{pn 8}%
\special{pa 4744 3760}%
\special{pa 4626 3634}%
\special{pa 4520 3516}%
\special{pa 4420 3406}%
\special{pa 4330 3304}%
\special{pa 4248 3210}%
\special{pa 4174 3120}%
\special{pa 4106 3038}%
\special{pa 4044 2960}%
\special{pa 3990 2886}%
\special{pa 3940 2818}%
\special{pa 3896 2752}%
\special{pa 3856 2690}%
\special{pa 3822 2630}%
\special{pa 3792 2574}%
\special{pa 3766 2520}%
\special{pa 3746 2468}%
\special{pa 3728 2416}%
\special{pa 3716 2368}%
\special{pa 3706 2318}%
\special{pa 3702 2270}%
\special{pa 3700 2222}%
\special{pa 3704 2174}%
\special{pa 3710 2126}%
\special{pa 3720 2076}%
\special{pa 3734 2026}%
\special{pa 3752 1974}%
\special{pa 3776 1922}%
\special{pa 3802 1866}%
\special{pa 3834 1810}%
\special{pa 3870 1750}%
\special{pa 3910 1686}%
\special{pa 3956 1620}%
\special{pa 4008 1548}%
\special{pa 4066 1474}%
\special{pa 4130 1394}%
\special{pa 4200 1310}%
\special{pa 4278 1218}%
\special{pa 4362 1120}%
\special{pa 4456 1016}%
\special{pa 4556 904}%
\special{pa 4668 782}%
\special{pa 4668 782}%
\special{pa 4670 780}%
\special{pa 4672 778}%
\special{pa 4672 778}%
\special{pa 4674 776}%
\special{pa 4674 776}%
\special{pa 4676 774}%
\special{pa 4676 772}%
\special{pa 4678 772}%
\special{pa 4680 770}%
\special{pa 4680 768}%
\special{pa 4682 768}%
\special{pa 4682 766}%
\special{pa 4684 766}%
\special{pa 4686 764}%
\special{pa 4686 762}%
\special{pa 4688 762}%
\special{pa 4688 760}%
\special{pa 4690 758}%
\special{pa 4692 758}%
\special{pa 4692 756}%
\special{pa 4694 754}%
\special{pa 4694 754}%
\special{pa 4696 752}%
\special{pa 4696 752}%
\special{pa 4698 750}%
\special{pa 4700 748}%
\special{pa 4700 748}%
\special{pa 4702 746}%
\special{pa 4702 744}%
\special{pa 4704 744}%
\special{pa 4706 742}%
\special{pa 4706 740}%
\special{pa 4708 740}%
\special{pa 4708 738}%
\special{pa 4710 738}%
\special{pa 4712 736}%
\special{pa 4712 734}%
\special{pa 4714 734}%
\special{pa 4714 732}%
\special{pa 4716 730}%
\special{pa 4718 730}%
\special{pa 4718 728}%
\special{pa 4720 726}%
\special{pa 4720 726}%
\special{pa 4722 724}%
\special{pa 4724 722}%
\special{pa 4724 722}%
\special{pa 4726 720}%
\special{pa 4726 718}%
\special{pa 4728 718}%
\special{pa 4730 716}%
\special{pa 4730 714}%
\special{pa 4732 714}%
\special{pa 4732 712}%
\special{pa 4734 712}%
\special{pa 4736 710}%
\special{pa 4736 708}%
\special{pa 4738 708}%
\special{pa 4738 706}%
\special{pa 4740 704}%
\special{pa 4742 704}%
\special{pa 4742 702}%
\special{pa 4744 700}%
\special{pa 4746 700}%
\special{pa 4746 698}%
\special{pa 4748 696}%
\special{pa 4748 696}%
\special{pa 4750 694}%
\special{pa 4752 692}%
\special{pa 4752 692}%
\special{pa 4754 690}%
\special{pa 4754 688}%
\special{pa 4756 688}%
\special{pa 4758 686}%
\special{pa 4758 684}%
\special{pa 4760 684}%
\special{pa 4762 682}%
\special{pa 4762 680}%
\special{pa 4764 680}%
\special{pa 4764 678}%
\special{pa 4766 676}%
\special{pa 4768 676}%
\special{pa 4768 674}%
\special{pa 4770 672}%
\special{pa 4772 672}%
\special{pa 4772 670}%
\special{pa 4774 668}%
\special{pa 4774 668}%
\special{pa 4776 666}%
\special{pa 4778 664}%
\special{pa 4778 664}%
\special{pa 4780 662}%
\special{sp}%
% FUNC 2 0 3 0
% 10 1380 580 4780 3760 3100 2230 3700 2230 3100 1630 1380 580 4780 3760 50 3 0 2 0 0
% -(exp(t)+exp(-t))/2///(exp(t)-exp(-t))/2///-10///10
\special{pn 8}%
\special{pa 1458 3760}%
\special{pa 1574 3634}%
\special{pa 1682 3516}%
\special{pa 1780 3406}%
\special{pa 1870 3304}%
\special{pa 1952 3210}%
\special{pa 2028 3120}%
\special{pa 2094 3038}%
\special{pa 2156 2960}%
\special{pa 2212 2886}%
\special{pa 2262 2818}%
\special{pa 2306 2752}%
\special{pa 2344 2690}%
\special{pa 2380 2630}%
\special{pa 2408 2574}%
\special{pa 2434 2520}%
\special{pa 2456 2468}%
\special{pa 2472 2416}%
\special{pa 2486 2368}%
\special{pa 2494 2318}%
\special{pa 2500 2270}%
\special{pa 2500 2222}%
\special{pa 2498 2174}%
\special{pa 2492 2126}%
\special{pa 2480 2076}%
\special{pa 2466 2026}%
\special{pa 2448 1974}%
\special{pa 2426 1922}%
\special{pa 2398 1866}%
\special{pa 2368 1810}%
\special{pa 2332 1750}%
\special{pa 2290 1686}%
\special{pa 2244 1620}%
\special{pa 2192 1548}%
\special{pa 2134 1474}%
\special{pa 2072 1394}%
\special{pa 2000 1310}%
\special{pa 1924 1218}%
\special{pa 1838 1120}%
\special{pa 1746 1016}%
\special{pa 1644 904}%
\special{pa 1534 782}%
\special{pa 1412 652}%
\special{pa 1412 652}%
\special{pa 1410 650}%
\special{pa 1408 648}%
\special{pa 1408 648}%
\special{pa 1406 646}%
\special{pa 1404 644}%
\special{pa 1404 644}%
\special{pa 1402 642}%
\special{pa 1400 640}%
\special{pa 1400 640}%
\special{pa 1398 638}%
\special{pa 1398 636}%
\special{pa 1396 634}%
\special{pa 1394 634}%
\special{pa 1394 632}%
\special{pa 1392 630}%
\special{pa 1390 630}%
\special{pa 1390 628}%
\special{pa 1388 626}%
\special{pa 1386 626}%
\special{pa 1386 624}%
\special{pa 1384 622}%
\special{pa 1384 622}%
\special{pa 1382 620}%
\special{pa 1380 618}%
\special{sp}%
% FUNC 2 2 3 0
% 10 1380 580 4780 3760 3100 2230 3700 2230 3100 1630 1380 580 4780 3760 0 1 0 0 0 0
% x
\special{pn 8}%
\special{pa 1570 3760}%
\special{pa 1570 3760}%
\special{pa 1576 3756}%
\special{pa 1580 3750}%
\special{pa 1586 3746}%
\special{pa 1590 3740}%
\special{pa 1596 3736}%
\special{pa 1600 3730}%
\special{pa 1606 3726}%
\special{pa 1610 3720}%
\special{pa 1616 3716}%
\special{pa 1620 3710}%
\special{pa 1626 3706}%
\special{pa 1630 3700}%
\special{pa 1636 3696}%
\special{pa 1640 3690}%
\special{pa 1646 3686}%
\special{pa 1650 3680}%
\special{pa 1656 3676}%
\special{pa 1660 3670}%
\special{pa 1666 3666}%
\special{pa 1670 3660}%
\special{pa 1676 3656}%
\special{pa 1680 3650}%
\special{pa 1686 3646}%
\special{pa 1690 3640}%
\special{pa 1696 3636}%
\special{pa 1700 3630}%
\special{pa 1706 3626}%
\special{pa 1710 3620}%
\special{pa 1716 3616}%
\special{pa 1720 3610}%
\special{pa 1726 3606}%
\special{pa 1730 3600}%
\special{pa 1736 3596}%
\special{pa 1740 3590}%
\special{pa 1746 3586}%
\special{pa 1750 3580}%
\special{pa 1756 3576}%
\special{pa 1760 3570}%
\special{pa 1766 3566}%
\special{pa 1770 3560}%
\special{pa 1776 3556}%
\special{pa 1780 3550}%
\special{pa 1786 3546}%
\special{pa 1790 3540}%
\special{pa 1796 3536}%
\special{pa 1800 3530}%
\special{pa 1806 3526}%
\special{pa 1810 3520}%
\special{pa 1816 3516}%
\special{pa 1820 3510}%
\special{pa 1826 3506}%
\special{pa 1830 3500}%
\special{pa 1836 3496}%
\special{pa 1840 3490}%
\special{pa 1846 3486}%
\special{pa 1850 3480}%
\special{pa 1856 3476}%
\special{pa 1860 3470}%
\special{pa 1866 3466}%
\special{pa 1870 3460}%
\special{pa 1876 3456}%
\special{pa 1880 3450}%
\special{pa 1886 3446}%
\special{pa 1890 3440}%
\special{pa 1896 3436}%
\special{pa 1900 3430}%
\special{pa 1906 3426}%
\special{pa 1910 3420}%
\special{pa 1916 3416}%
\special{pa 1920 3410}%
\special{pa 1926 3406}%
\special{pa 1930 3400}%
\special{pa 1936 3396}%
\special{pa 1940 3390}%
\special{pa 1946 3386}%
\special{pa 1950 3380}%
\special{pa 1956 3376}%
\special{pa 1960 3370}%
\special{pa 1966 3366}%
\special{pa 1970 3360}%
\special{pa 1976 3356}%
\special{pa 1980 3350}%
\special{pa 1986 3346}%
\special{pa 1990 3340}%
\special{pa 1996 3336}%
\special{pa 2000 3330}%
\special{pa 2006 3326}%
\special{pa 2010 3320}%
\special{pa 2016 3316}%
\special{pa 2020 3310}%
\special{pa 2026 3306}%
\special{pa 2030 3300}%
\special{pa 2036 3296}%
\special{pa 2040 3290}%
\special{pa 2046 3286}%
\special{pa 2050 3280}%
\special{pa 2056 3276}%
\special{pa 2060 3270}%
\special{pa 2066 3266}%
\special{pa 2070 3260}%
\special{pa 2076 3256}%
\special{pa 2080 3250}%
\special{pa 2086 3246}%
\special{pa 2090 3240}%
\special{pa 2096 3236}%
\special{pa 2100 3230}%
\special{pa 2106 3226}%
\special{pa 2110 3220}%
\special{pa 2116 3216}%
\special{pa 2120 3210}%
\special{pa 2126 3206}%
\special{pa 2130 3200}%
\special{pa 2136 3196}%
\special{pa 2140 3190}%
\special{pa 2146 3186}%
\special{pa 2150 3180}%
\special{pa 2156 3176}%
\special{pa 2160 3170}%
\special{pa 2166 3166}%
\special{pa 2170 3160}%
\special{pa 2176 3156}%
\special{pa 2180 3150}%
\special{pa 2186 3146}%
\special{pa 2190 3140}%
\special{pa 2196 3136}%
\special{pa 2200 3130}%
\special{pa 2206 3126}%
\special{pa 2210 3120}%
\special{pa 2216 3116}%
\special{pa 2220 3110}%
\special{pa 2226 3106}%
\special{pa 2230 3100}%
\special{pa 2236 3096}%
\special{pa 2240 3090}%
\special{pa 2246 3086}%
\special{pa 2250 3080}%
\special{pa 2256 3076}%
\special{pa 2260 3070}%
\special{pa 2266 3066}%
\special{pa 2270 3060}%
\special{pa 2276 3056}%
\special{pa 2280 3050}%
\special{pa 2286 3046}%
\special{pa 2290 3040}%
\special{pa 2296 3036}%
\special{pa 2300 3030}%
\special{pa 2306 3026}%
\special{pa 2310 3020}%
\special{pa 2316 3016}%
\special{pa 2320 3010}%
\special{pa 2326 3006}%
\special{pa 2330 3000}%
\special{pa 2336 2996}%
\special{pa 2340 2990}%
\special{pa 2346 2986}%
\special{pa 2350 2980}%
\special{pa 2356 2976}%
\special{pa 2360 2970}%
\special{pa 2366 2966}%
\special{pa 2370 2960}%
\special{pa 2376 2956}%
\special{pa 2380 2950}%
\special{pa 2386 2946}%
\special{pa 2390 2940}%
\special{pa 2396 2936}%
\special{pa 2400 2930}%
\special{pa 2406 2926}%
\special{pa 2410 2920}%
\special{pa 2416 2916}%
\special{pa 2420 2910}%
\special{pa 2426 2906}%
\special{pa 2430 2900}%
\special{pa 2436 2896}%
\special{pa 2440 2890}%
\special{pa 2446 2886}%
\special{pa 2450 2880}%
\special{pa 2456 2876}%
\special{pa 2460 2870}%
\special{pa 2466 2866}%
\special{pa 2470 2860}%
\special{pa 2476 2856}%
\special{pa 2480 2850}%
\special{pa 2486 2846}%
\special{pa 2490 2840}%
\special{pa 2496 2836}%
\special{pa 2500 2830}%
\special{pa 2506 2826}%
\special{pa 2510 2820}%
\special{pa 2516 2816}%
\special{pa 2520 2810}%
\special{pa 2526 2806}%
\special{pa 2530 2800}%
\special{pa 2536 2796}%
\special{pa 2540 2790}%
\special{pa 2546 2786}%
\special{pa 2550 2780}%
\special{pa 2556 2776}%
\special{pa 2560 2770}%
\special{pa 2566 2766}%
\special{pa 2570 2760}%
\special{pa 2576 2756}%
\special{pa 2580 2750}%
\special{pa 2586 2746}%
\special{pa 2590 2740}%
\special{pa 2596 2736}%
\special{pa 2600 2730}%
\special{pa 2606 2726}%
\special{pa 2610 2720}%
\special{pa 2616 2716}%
\special{pa 2620 2710}%
\special{pa 2626 2706}%
\special{pa 2630 2700}%
\special{pa 2636 2696}%
\special{pa 2640 2690}%
\special{pa 2646 2686}%
\special{pa 2650 2680}%
\special{pa 2656 2676}%
\special{pa 2660 2670}%
\special{pa 2666 2666}%
\special{pa 2670 2660}%
\special{pa 2676 2656}%
\special{pa 2680 2650}%
\special{pa 2686 2646}%
\special{pa 2690 2640}%
\special{pa 2696 2636}%
\special{pa 2700 2630}%
\special{pa 2706 2626}%
\special{pa 2710 2620}%
\special{pa 2716 2616}%
\special{pa 2720 2610}%
\special{pa 2726 2606}%
\special{pa 2730 2600}%
\special{pa 2736 2596}%
\special{pa 2740 2590}%
\special{pa 2746 2586}%
\special{pa 2750 2580}%
\special{pa 2756 2576}%
\special{pa 2760 2570}%
\special{pa 2766 2566}%
\special{pa 2770 2560}%
\special{pa 2776 2556}%
\special{pa 2780 2550}%
\special{pa 2786 2546}%
\special{pa 2790 2540}%
\special{pa 2796 2536}%
\special{pa 2800 2530}%
\special{pa 2806 2526}%
\special{pa 2810 2520}%
\special{pa 2816 2516}%
\special{pa 2820 2510}%
\special{pa 2826 2506}%
\special{pa 2830 2500}%
\special{pa 2836 2496}%
\special{pa 2840 2490}%
\special{pa 2846 2486}%
\special{pa 2850 2480}%
\special{pa 2856 2476}%
\special{pa 2860 2470}%
\special{pa 2866 2466}%
\special{pa 2870 2460}%
\special{pa 2876 2456}%
\special{pa 2880 2450}%
\special{pa 2886 2446}%
\special{pa 2890 2440}%
\special{pa 2896 2436}%
\special{pa 2900 2430}%
\special{pa 2906 2426}%
\special{pa 2910 2420}%
\special{pa 2916 2416}%
\special{pa 2920 2410}%
\special{pa 2926 2406}%
\special{pa 2930 2400}%
\special{pa 2936 2396}%
\special{pa 2940 2390}%
\special{pa 2946 2386}%
\special{pa 2950 2380}%
\special{pa 2956 2376}%
\special{pa 2960 2370}%
\special{pa 2966 2366}%
\special{pa 2970 2360}%
\special{pa 2976 2356}%
\special{pa 2980 2350}%
\special{pa 2986 2346}%
\special{pa 2990 2340}%
\special{pa 2996 2336}%
\special{pa 3000 2330}%
\special{pa 3006 2326}%
\special{pa 3010 2320}%
\special{pa 3016 2316}%
\special{pa 3020 2310}%
\special{pa 3026 2306}%
\special{pa 3030 2300}%
\special{pa 3036 2296}%
\special{pa 3040 2290}%
\special{pa 3046 2286}%
\special{pa 3050 2280}%
\special{pa 3056 2276}%
\special{pa 3060 2270}%
\special{pa 3066 2266}%
\special{pa 3070 2260}%
\special{pa 3076 2256}%
\special{pa 3080 2250}%
\special{pa 3086 2246}%
\special{pa 3090 2240}%
\special{pa 3096 2236}%
\special{pa 3100 2230}%
\special{pa 3106 2226}%
\special{pa 3110 2220}%
\special{pa 3116 2216}%
\special{pa 3120 2210}%
\special{pa 3126 2206}%
\special{pa 3130 2200}%
\special{pa 3136 2196}%
\special{pa 3140 2190}%
\special{pa 3146 2186}%
\special{pa 3150 2180}%
\special{pa 3156 2176}%
\special{pa 3160 2170}%
\special{pa 3166 2166}%
\special{pa 3170 2160}%
\special{pa 3176 2156}%
\special{pa 3180 2150}%
\special{pa 3186 2146}%
\special{pa 3190 2140}%
\special{pa 3196 2136}%
\special{pa 3200 2130}%
\special{pa 3206 2126}%
\special{pa 3210 2120}%
\special{pa 3216 2116}%
\special{pa 3220 2110}%
\special{pa 3226 2106}%
\special{pa 3230 2100}%
\special{pa 3236 2096}%
\special{pa 3240 2090}%
\special{pa 3246 2086}%
\special{pa 3250 2080}%
\special{pa 3256 2076}%
\special{pa 3260 2070}%
\special{pa 3266 2066}%
\special{pa 3270 2060}%
\special{pa 3276 2056}%
\special{pa 3280 2050}%
\special{pa 3286 2046}%
\special{pa 3290 2040}%
\special{pa 3296 2036}%
\special{pa 3300 2030}%
\special{pa 3306 2026}%
\special{pa 3310 2020}%
\special{pa 3316 2016}%
\special{pa 3320 2010}%
\special{pa 3326 2006}%
\special{pa 3330 2000}%
\special{pa 3336 1996}%
\special{pa 3340 1990}%
\special{pa 3346 1986}%
\special{pa 3350 1980}%
\special{pa 3356 1976}%
\special{pa 3360 1970}%
\special{pa 3366 1966}%
\special{pa 3370 1960}%
\special{pa 3376 1956}%
\special{pa 3380 1950}%
\special{pa 3386 1946}%
\special{pa 3390 1940}%
\special{pa 3396 1936}%
\special{pa 3400 1930}%
\special{pa 3406 1926}%
\special{pa 3410 1920}%
\special{pa 3416 1916}%
\special{pa 3420 1910}%
\special{pa 3426 1906}%
\special{pa 3430 1900}%
\special{pa 3436 1896}%
\special{pa 3440 1890}%
\special{pa 3446 1886}%
\special{pa 3450 1880}%
\special{pa 3456 1876}%
\special{pa 3460 1870}%
\special{pa 3466 1866}%
\special{pa 3470 1860}%
\special{pa 3476 1856}%
\special{pa 3480 1850}%
\special{pa 3486 1846}%
\special{pa 3490 1840}%
\special{pa 3496 1836}%
\special{pa 3500 1830}%
\special{pa 3506 1826}%
\special{pa 3510 1820}%
\special{pa 3516 1816}%
\special{pa 3520 1810}%
\special{pa 3526 1806}%
\special{pa 3530 1800}%
\special{pa 3536 1796}%
\special{pa 3540 1790}%
\special{pa 3546 1786}%
\special{pa 3550 1780}%
\special{pa 3556 1776}%
\special{pa 3560 1770}%
\special{pa 3566 1766}%
\special{pa 3570 1760}%
\special{pa 3576 1756}%
\special{pa 3580 1750}%
\special{pa 3586 1746}%
\special{pa 3590 1740}%
\special{pa 3596 1736}%
\special{pa 3600 1730}%
\special{pa 3606 1726}%
\special{pa 3610 1720}%
\special{pa 3616 1716}%
\special{pa 3620 1710}%
\special{pa 3626 1706}%
\special{pa 3630 1700}%
\special{pa 3636 1696}%
\special{pa 3640 1690}%
\special{pa 3646 1686}%
\special{pa 3650 1680}%
\special{pa 3656 1676}%
\special{pa 3660 1670}%
\special{pa 3666 1666}%
\special{pa 3670 1660}%
\special{pa 3676 1656}%
\special{pa 3680 1650}%
\special{pa 3686 1646}%
\special{pa 3690 1640}%
\special{pa 3696 1636}%
\special{pa 3700 1630}%
\special{pa 3706 1626}%
\special{pa 3710 1620}%
\special{pa 3716 1616}%
\special{pa 3720 1610}%
\special{pa 3726 1606}%
\special{pa 3730 1600}%
\special{pa 3736 1596}%
\special{pa 3740 1590}%
\special{pa 3746 1586}%
\special{pa 3750 1580}%
\special{pa 3756 1576}%
\special{pa 3760 1570}%
\special{pa 3766 1566}%
\special{pa 3770 1560}%
\special{pa 3776 1556}%
\special{pa 3780 1550}%
\special{pa 3786 1546}%
\special{pa 3790 1540}%
\special{pa 3796 1536}%
\special{pa 3800 1530}%
\special{pa 3806 1526}%
\special{pa 3810 1520}%
\special{pa 3816 1516}%
\special{pa 3820 1510}%
\special{pa 3826 1506}%
\special{pa 3830 1500}%
\special{pa 3836 1496}%
\special{pa 3840 1490}%
\special{pa 3846 1486}%
\special{pa 3850 1480}%
\special{pa 3856 1476}%
\special{pa 3860 1470}%
\special{pa 3866 1466}%
\special{pa 3870 1460}%
\special{pa 3876 1456}%
\special{pa 3880 1450}%
\special{pa 3886 1446}%
\special{pa 3890 1440}%
\special{pa 3896 1436}%
\special{pa 3900 1430}%
\special{pa 3906 1426}%
\special{pa 3910 1420}%
\special{pa 3916 1416}%
\special{pa 3920 1410}%
\special{pa 3926 1406}%
\special{pa 3930 1400}%
\special{pa 3936 1396}%
\special{pa 3940 1390}%
\special{pa 3946 1386}%
\special{pa 3950 1380}%
\special{pa 3956 1376}%
\special{pa 3960 1370}%
\special{pa 3966 1366}%
\special{pa 3970 1360}%
\special{pa 3976 1356}%
\special{pa 3980 1350}%
\special{pa 3986 1346}%
\special{pa 3990 1340}%
\special{pa 3996 1336}%
\special{pa 4000 1330}%
\special{pa 4006 1326}%
\special{pa 4010 1320}%
\special{pa 4016 1316}%
\special{pa 4020 1310}%
\special{pa 4026 1306}%
\special{pa 4030 1300}%
\special{pa 4036 1296}%
\special{pa 4040 1290}%
\special{pa 4046 1286}%
\special{pa 4050 1280}%
\special{pa 4056 1276}%
\special{pa 4060 1270}%
\special{pa 4066 1266}%
\special{pa 4070 1260}%
\special{pa 4076 1256}%
\special{pa 4080 1250}%
\special{pa 4086 1246}%
\special{pa 4090 1240}%
\special{pa 4096 1236}%
\special{pa 4100 1230}%
\special{pa 4106 1226}%
\special{pa 4110 1220}%
\special{pa 4116 1216}%
\special{pa 4120 1210}%
\special{pa 4126 1206}%
\special{pa 4130 1200}%
\special{pa 4136 1196}%
\special{pa 4140 1190}%
\special{pa 4146 1186}%
\special{pa 4150 1180}%
\special{pa 4156 1176}%
\special{pa 4160 1170}%
\special{pa 4166 1166}%
\special{pa 4170 1160}%
\special{pa 4176 1156}%
\special{pa 4180 1150}%
\special{pa 4186 1146}%
\special{pa 4190 1140}%
\special{pa 4196 1136}%
\special{pa 4200 1130}%
\special{pa 4206 1126}%
\special{pa 4210 1120}%
\special{pa 4216 1116}%
\special{pa 4220 1110}%
\special{pa 4226 1106}%
\special{pa 4230 1100}%
\special{pa 4236 1096}%
\special{pa 4240 1090}%
\special{pa 4246 1086}%
\special{pa 4250 1080}%
\special{pa 4256 1076}%
\special{pa 4260 1070}%
\special{pa 4266 1066}%
\special{pa 4270 1060}%
\special{pa 4276 1056}%
\special{pa 4280 1050}%
\special{pa 4286 1046}%
\special{pa 4290 1040}%
\special{pa 4296 1036}%
\special{pa 4300 1030}%
\special{pa 4306 1026}%
\special{pa 4310 1020}%
\special{pa 4316 1016}%
\special{pa 4320 1010}%
\special{pa 4326 1006}%
\special{pa 4330 1000}%
\special{pa 4336 996}%
\special{pa 4340 990}%
\special{pa 4346 986}%
\special{pa 4350 980}%
\special{pa 4356 976}%
\special{pa 4360 970}%
\special{pa 4366 966}%
\special{pa 4370 960}%
\special{pa 4376 956}%
\special{pa 4380 950}%
\special{pa 4386 946}%
\special{pa 4390 940}%
\special{pa 4396 936}%
\special{pa 4400 930}%
\special{pa 4406 926}%
\special{pa 4410 920}%
\special{pa 4416 916}%
\special{pa 4420 910}%
\special{pa 4426 906}%
\special{pa 4430 900}%
\special{pa 4436 896}%
\special{pa 4440 890}%
\special{pa 4446 886}%
\special{pa 4450 880}%
\special{pa 4456 876}%
\special{pa 4460 870}%
\special{pa 4466 866}%
\special{pa 4470 860}%
\special{pa 4476 856}%
\special{pa 4480 850}%
\special{pa 4486 846}%
\special{pa 4490 840}%
\special{pa 4496 836}%
\special{pa 4500 830}%
\special{pa 4506 826}%
\special{pa 4510 820}%
\special{pa 4516 816}%
\special{pa 4520 810}%
\special{pa 4526 806}%
\special{pa 4530 800}%
\special{pa 4536 796}%
\special{pa 4540 790}%
\special{pa 4546 786}%
\special{pa 4550 780}%
\special{pa 4556 776}%
\special{pa 4560 770}%
\special{pa 4566 766}%
\special{pa 4570 760}%
\special{pa 4576 756}%
\special{pa 4580 750}%
\special{pa 4586 746}%
\special{pa 4590 740}%
\special{pa 4596 736}%
\special{pa 4600 730}%
\special{pa 4606 726}%
\special{pa 4610 720}%
\special{pa 4616 716}%
\special{pa 4620 710}%
\special{pa 4626 706}%
\special{pa 4630 700}%
\special{pa 4636 696}%
\special{pa 4640 690}%
\special{pa 4646 686}%
\special{pa 4650 680}%
\special{pa 4656 676}%
\special{pa 4660 670}%
\special{pa 4666 666}%
\special{pa 4670 660}%
\special{pa 4676 656}%
\special{pa 4680 650}%
\special{pa 4686 646}%
\special{pa 4690 640}%
\special{pa 4696 636}%
\special{pa 4700 630}%
\special{pa 4706 626}%
\special{pa 4710 620}%
\special{pa 4716 616}%
\special{pa 4720 610}%
\special{pa 4726 606}%
\special{pa 4730 600}%
\special{pa 4736 596}%
\special{pa 4740 590}%
\special{pa 4746 586}%
\special{pa 4750 580}%
\special{pa 4750 580}%
\special{sp -0.045}%
% FUNC 2 2 3 0
% 10 1380 580 4780 3760 3100 2230 3700 2230 3100 1630 1380 580 4780 3760 0 1 0 0 0 0
% -x
\special{pn 8}%
\special{pa 1450 580}%
\special{pa 1450 580}%
\special{pa 1456 586}%
\special{pa 1460 590}%
\special{pa 1466 596}%
\special{pa 1470 600}%
\special{pa 1476 606}%
\special{pa 1480 610}%
\special{pa 1486 616}%
\special{pa 1490 620}%
\special{pa 1496 626}%
\special{pa 1500 630}%
\special{pa 1506 636}%
\special{pa 1510 640}%
\special{pa 1516 646}%
\special{pa 1520 650}%
\special{pa 1526 656}%
\special{pa 1530 660}%
\special{pa 1536 666}%
\special{pa 1540 670}%
\special{pa 1546 676}%
\special{pa 1550 680}%
\special{pa 1556 686}%
\special{pa 1560 690}%
\special{pa 1566 696}%
\special{pa 1570 700}%
\special{pa 1576 706}%
\special{pa 1580 710}%
\special{pa 1586 716}%
\special{pa 1590 720}%
\special{pa 1596 726}%
\special{pa 1600 730}%
\special{pa 1606 736}%
\special{pa 1610 740}%
\special{pa 1616 746}%
\special{pa 1620 750}%
\special{pa 1626 756}%
\special{pa 1630 760}%
\special{pa 1636 766}%
\special{pa 1640 770}%
\special{pa 1646 776}%
\special{pa 1650 780}%
\special{pa 1656 786}%
\special{pa 1660 790}%
\special{pa 1666 796}%
\special{pa 1670 800}%
\special{pa 1676 806}%
\special{pa 1680 810}%
\special{pa 1686 816}%
\special{pa 1690 820}%
\special{pa 1696 826}%
\special{pa 1700 830}%
\special{pa 1706 836}%
\special{pa 1710 840}%
\special{pa 1716 846}%
\special{pa 1720 850}%
\special{pa 1726 856}%
\special{pa 1730 860}%
\special{pa 1736 866}%
\special{pa 1740 870}%
\special{pa 1746 876}%
\special{pa 1750 880}%
\special{pa 1756 886}%
\special{pa 1760 890}%
\special{pa 1766 896}%
\special{pa 1770 900}%
\special{pa 1776 906}%
\special{pa 1780 910}%
\special{pa 1786 916}%
\special{pa 1790 920}%
\special{pa 1796 926}%
\special{pa 1800 930}%
\special{pa 1806 936}%
\special{pa 1810 940}%
\special{pa 1816 946}%
\special{pa 1820 950}%
\special{pa 1826 956}%
\special{pa 1830 960}%
\special{pa 1836 966}%
\special{pa 1840 970}%
\special{pa 1846 976}%
\special{pa 1850 980}%
\special{pa 1856 986}%
\special{pa 1860 990}%
\special{pa 1866 996}%
\special{pa 1870 1000}%
\special{pa 1876 1006}%
\special{pa 1880 1010}%
\special{pa 1886 1016}%
\special{pa 1890 1020}%
\special{pa 1896 1026}%
\special{pa 1900 1030}%
\special{pa 1906 1036}%
\special{pa 1910 1040}%
\special{pa 1916 1046}%
\special{pa 1920 1050}%
\special{pa 1926 1056}%
\special{pa 1930 1060}%
\special{pa 1936 1066}%
\special{pa 1940 1070}%
\special{pa 1946 1076}%
\special{pa 1950 1080}%
\special{pa 1956 1086}%
\special{pa 1960 1090}%
\special{pa 1966 1096}%
\special{pa 1970 1100}%
\special{pa 1976 1106}%
\special{pa 1980 1110}%
\special{pa 1986 1116}%
\special{pa 1990 1120}%
\special{pa 1996 1126}%
\special{pa 2000 1130}%
\special{pa 2006 1136}%
\special{pa 2010 1140}%
\special{pa 2016 1146}%
\special{pa 2020 1150}%
\special{pa 2026 1156}%
\special{pa 2030 1160}%
\special{pa 2036 1166}%
\special{pa 2040 1170}%
\special{pa 2046 1176}%
\special{pa 2050 1180}%
\special{pa 2056 1186}%
\special{pa 2060 1190}%
\special{pa 2066 1196}%
\special{pa 2070 1200}%
\special{pa 2076 1206}%
\special{pa 2080 1210}%
\special{pa 2086 1216}%
\special{pa 2090 1220}%
\special{pa 2096 1226}%
\special{pa 2100 1230}%
\special{pa 2106 1236}%
\special{pa 2110 1240}%
\special{pa 2116 1246}%
\special{pa 2120 1250}%
\special{pa 2126 1256}%
\special{pa 2130 1260}%
\special{pa 2136 1266}%
\special{pa 2140 1270}%
\special{pa 2146 1276}%
\special{pa 2150 1280}%
\special{pa 2156 1286}%
\special{pa 2160 1290}%
\special{pa 2166 1296}%
\special{pa 2170 1300}%
\special{pa 2176 1306}%
\special{pa 2180 1310}%
\special{pa 2186 1316}%
\special{pa 2190 1320}%
\special{pa 2196 1326}%
\special{pa 2200 1330}%
\special{pa 2206 1336}%
\special{pa 2210 1340}%
\special{pa 2216 1346}%
\special{pa 2220 1350}%
\special{pa 2226 1356}%
\special{pa 2230 1360}%
\special{pa 2236 1366}%
\special{pa 2240 1370}%
\special{pa 2246 1376}%
\special{pa 2250 1380}%
\special{pa 2256 1386}%
\special{pa 2260 1390}%
\special{pa 2266 1396}%
\special{pa 2270 1400}%
\special{pa 2276 1406}%
\special{pa 2280 1410}%
\special{pa 2286 1416}%
\special{pa 2290 1420}%
\special{pa 2296 1426}%
\special{pa 2300 1430}%
\special{pa 2306 1436}%
\special{pa 2310 1440}%
\special{pa 2316 1446}%
\special{pa 2320 1450}%
\special{pa 2326 1456}%
\special{pa 2330 1460}%
\special{pa 2336 1466}%
\special{pa 2340 1470}%
\special{pa 2346 1476}%
\special{pa 2350 1480}%
\special{pa 2356 1486}%
\special{pa 2360 1490}%
\special{pa 2366 1496}%
\special{pa 2370 1500}%
\special{pa 2376 1506}%
\special{pa 2380 1510}%
\special{pa 2386 1516}%
\special{pa 2390 1520}%
\special{pa 2396 1526}%
\special{pa 2400 1530}%
\special{pa 2406 1536}%
\special{pa 2410 1540}%
\special{pa 2416 1546}%
\special{pa 2420 1550}%
\special{pa 2426 1556}%
\special{pa 2430 1560}%
\special{pa 2436 1566}%
\special{pa 2440 1570}%
\special{pa 2446 1576}%
\special{pa 2450 1580}%
\special{pa 2456 1586}%
\special{pa 2460 1590}%
\special{pa 2466 1596}%
\special{pa 2470 1600}%
\special{pa 2476 1606}%
\special{pa 2480 1610}%
\special{pa 2486 1616}%
\special{pa 2490 1620}%
\special{pa 2496 1626}%
\special{pa 2500 1630}%
\special{pa 2506 1636}%
\special{pa 2510 1640}%
\special{pa 2516 1646}%
\special{pa 2520 1650}%
\special{pa 2526 1656}%
\special{pa 2530 1660}%
\special{pa 2536 1666}%
\special{pa 2540 1670}%
\special{pa 2546 1676}%
\special{pa 2550 1680}%
\special{pa 2556 1686}%
\special{pa 2560 1690}%
\special{pa 2566 1696}%
\special{pa 2570 1700}%
\special{pa 2576 1706}%
\special{pa 2580 1710}%
\special{pa 2586 1716}%
\special{pa 2590 1720}%
\special{pa 2596 1726}%
\special{pa 2600 1730}%
\special{pa 2606 1736}%
\special{pa 2610 1740}%
\special{pa 2616 1746}%
\special{pa 2620 1750}%
\special{pa 2626 1756}%
\special{pa 2630 1760}%
\special{pa 2636 1766}%
\special{pa 2640 1770}%
\special{pa 2646 1776}%
\special{pa 2650 1780}%
\special{pa 2656 1786}%
\special{pa 2660 1790}%
\special{pa 2666 1796}%
\special{pa 2670 1800}%
\special{pa 2676 1806}%
\special{pa 2680 1810}%
\special{pa 2686 1816}%
\special{pa 2690 1820}%
\special{pa 2696 1826}%
\special{pa 2700 1830}%
\special{pa 2706 1836}%
\special{pa 2710 1840}%
\special{pa 2716 1846}%
\special{pa 2720 1850}%
\special{pa 2726 1856}%
\special{pa 2730 1860}%
\special{pa 2736 1866}%
\special{pa 2740 1870}%
\special{pa 2746 1876}%
\special{pa 2750 1880}%
\special{pa 2756 1886}%
\special{pa 2760 1890}%
\special{pa 2766 1896}%
\special{pa 2770 1900}%
\special{pa 2776 1906}%
\special{pa 2780 1910}%
\special{pa 2786 1916}%
\special{pa 2790 1920}%
\special{pa 2796 1926}%
\special{pa 2800 1930}%
\special{pa 2806 1936}%
\special{pa 2810 1940}%
\special{pa 2816 1946}%
\special{pa 2820 1950}%
\special{pa 2826 1956}%
\special{pa 2830 1960}%
\special{pa 2836 1966}%
\special{pa 2840 1970}%
\special{pa 2846 1976}%
\special{pa 2850 1980}%
\special{pa 2856 1986}%
\special{pa 2860 1990}%
\special{pa 2866 1996}%
\special{pa 2870 2000}%
\special{pa 2876 2006}%
\special{pa 2880 2010}%
\special{pa 2886 2016}%
\special{pa 2890 2020}%
\special{pa 2896 2026}%
\special{pa 2900 2030}%
\special{pa 2906 2036}%
\special{pa 2910 2040}%
\special{pa 2916 2046}%
\special{pa 2920 2050}%
\special{pa 2926 2056}%
\special{pa 2930 2060}%
\special{pa 2936 2066}%
\special{pa 2940 2070}%
\special{pa 2946 2076}%
\special{pa 2950 2080}%
\special{pa 2956 2086}%
\special{pa 2960 2090}%
\special{pa 2966 2096}%
\special{pa 2970 2100}%
\special{pa 2976 2106}%
\special{pa 2980 2110}%
\special{pa 2986 2116}%
\special{pa 2990 2120}%
\special{pa 2996 2126}%
\special{pa 3000 2130}%
\special{pa 3006 2136}%
\special{pa 3010 2140}%
\special{pa 3016 2146}%
\special{pa 3020 2150}%
\special{pa 3026 2156}%
\special{pa 3030 2160}%
\special{pa 3036 2166}%
\special{pa 3040 2170}%
\special{pa 3046 2176}%
\special{pa 3050 2180}%
\special{pa 3056 2186}%
\special{pa 3060 2190}%
\special{pa 3066 2196}%
\special{pa 3070 2200}%
\special{pa 3076 2206}%
\special{pa 3080 2210}%
\special{pa 3086 2216}%
\special{pa 3090 2220}%
\special{pa 3096 2226}%
\special{pa 3100 2230}%
\special{pa 3106 2236}%
\special{pa 3110 2240}%
\special{pa 3116 2246}%
\special{pa 3120 2250}%
\special{pa 3126 2256}%
\special{pa 3130 2260}%
\special{pa 3136 2266}%
\special{pa 3140 2270}%
\special{pa 3146 2276}%
\special{pa 3150 2280}%
\special{pa 3156 2286}%
\special{pa 3160 2290}%
\special{pa 3166 2296}%
\special{pa 3170 2300}%
\special{pa 3176 2306}%
\special{pa 3180 2310}%
\special{pa 3186 2316}%
\special{pa 3190 2320}%
\special{pa 3196 2326}%
\special{pa 3200 2330}%
\special{pa 3206 2336}%
\special{pa 3210 2340}%
\special{pa 3216 2346}%
\special{pa 3220 2350}%
\special{pa 3226 2356}%
\special{pa 3230 2360}%
\special{pa 3236 2366}%
\special{pa 3240 2370}%
\special{pa 3246 2376}%
\special{pa 3250 2380}%
\special{pa 3256 2386}%
\special{pa 3260 2390}%
\special{pa 3266 2396}%
\special{pa 3270 2400}%
\special{pa 3276 2406}%
\special{pa 3280 2410}%
\special{pa 3286 2416}%
\special{pa 3290 2420}%
\special{pa 3296 2426}%
\special{pa 3300 2430}%
\special{pa 3306 2436}%
\special{pa 3310 2440}%
\special{pa 3316 2446}%
\special{pa 3320 2450}%
\special{pa 3326 2456}%
\special{pa 3330 2460}%
\special{pa 3336 2466}%
\special{pa 3340 2470}%
\special{pa 3346 2476}%
\special{pa 3350 2480}%
\special{pa 3356 2486}%
\special{pa 3360 2490}%
\special{pa 3366 2496}%
\special{pa 3370 2500}%
\special{pa 3376 2506}%
\special{pa 3380 2510}%
\special{pa 3386 2516}%
\special{pa 3390 2520}%
\special{pa 3396 2526}%
\special{pa 3400 2530}%
\special{pa 3406 2536}%
\special{pa 3410 2540}%
\special{pa 3416 2546}%
\special{pa 3420 2550}%
\special{pa 3426 2556}%
\special{pa 3430 2560}%
\special{pa 3436 2566}%
\special{pa 3440 2570}%
\special{pa 3446 2576}%
\special{pa 3450 2580}%
\special{pa 3456 2586}%
\special{pa 3460 2590}%
\special{pa 3466 2596}%
\special{pa 3470 2600}%
\special{pa 3476 2606}%
\special{pa 3480 2610}%
\special{pa 3486 2616}%
\special{pa 3490 2620}%
\special{pa 3496 2626}%
\special{pa 3500 2630}%
\special{pa 3506 2636}%
\special{pa 3510 2640}%
\special{pa 3516 2646}%
\special{pa 3520 2650}%
\special{pa 3526 2656}%
\special{pa 3530 2660}%
\special{pa 3536 2666}%
\special{pa 3540 2670}%
\special{pa 3546 2676}%
\special{pa 3550 2680}%
\special{pa 3556 2686}%
\special{pa 3560 2690}%
\special{pa 3566 2696}%
\special{pa 3570 2700}%
\special{pa 3576 2706}%
\special{pa 3580 2710}%
\special{pa 3586 2716}%
\special{pa 3590 2720}%
\special{pa 3596 2726}%
\special{pa 3600 2730}%
\special{pa 3606 2736}%
\special{pa 3610 2740}%
\special{pa 3616 2746}%
\special{pa 3620 2750}%
\special{pa 3626 2756}%
\special{pa 3630 2760}%
\special{pa 3636 2766}%
\special{pa 3640 2770}%
\special{pa 3646 2776}%
\special{pa 3650 2780}%
\special{pa 3656 2786}%
\special{pa 3660 2790}%
\special{pa 3666 2796}%
\special{pa 3670 2800}%
\special{pa 3676 2806}%
\special{pa 3680 2810}%
\special{pa 3686 2816}%
\special{pa 3690 2820}%
\special{pa 3696 2826}%
\special{pa 3700 2830}%
\special{pa 3706 2836}%
\special{pa 3710 2840}%
\special{pa 3716 2846}%
\special{pa 3720 2850}%
\special{pa 3726 2856}%
\special{pa 3730 2860}%
\special{pa 3736 2866}%
\special{pa 3740 2870}%
\special{pa 3746 2876}%
\special{pa 3750 2880}%
\special{pa 3756 2886}%
\special{pa 3760 2890}%
\special{pa 3766 2896}%
\special{pa 3770 2900}%
\special{pa 3776 2906}%
\special{pa 3780 2910}%
\special{pa 3786 2916}%
\special{pa 3790 2920}%
\special{pa 3796 2926}%
\special{pa 3800 2930}%
\special{pa 3806 2936}%
\special{pa 3810 2940}%
\special{pa 3816 2946}%
\special{pa 3820 2950}%
\special{pa 3826 2956}%
\special{pa 3830 2960}%
\special{pa 3836 2966}%
\special{pa 3840 2970}%
\special{pa 3846 2976}%
\special{pa 3850 2980}%
\special{pa 3856 2986}%
\special{pa 3860 2990}%
\special{pa 3866 2996}%
\special{pa 3870 3000}%
\special{pa 3876 3006}%
\special{pa 3880 3010}%
\special{pa 3886 3016}%
\special{pa 3890 3020}%
\special{pa 3896 3026}%
\special{pa 3900 3030}%
\special{pa 3906 3036}%
\special{pa 3910 3040}%
\special{pa 3916 3046}%
\special{pa 3920 3050}%
\special{pa 3926 3056}%
\special{pa 3930 3060}%
\special{pa 3936 3066}%
\special{pa 3940 3070}%
\special{pa 3946 3076}%
\special{pa 3950 3080}%
\special{pa 3956 3086}%
\special{pa 3960 3090}%
\special{pa 3966 3096}%
\special{pa 3970 3100}%
\special{pa 3976 3106}%
\special{pa 3980 3110}%
\special{pa 3986 3116}%
\special{pa 3990 3120}%
\special{pa 3996 3126}%
\special{pa 4000 3130}%
\special{pa 4006 3136}%
\special{pa 4010 3140}%
\special{pa 4016 3146}%
\special{pa 4020 3150}%
\special{pa 4026 3156}%
\special{pa 4030 3160}%
\special{pa 4036 3166}%
\special{pa 4040 3170}%
\special{pa 4046 3176}%
\special{pa 4050 3180}%
\special{pa 4056 3186}%
\special{pa 4060 3190}%
\special{pa 4066 3196}%
\special{pa 4070 3200}%
\special{pa 4076 3206}%
\special{pa 4080 3210}%
\special{pa 4086 3216}%
\special{pa 4090 3220}%
\special{pa 4096 3226}%
\special{pa 4100 3230}%
\special{pa 4106 3236}%
\special{pa 4110 3240}%
\special{pa 4116 3246}%
\special{pa 4120 3250}%
\special{pa 4126 3256}%
\special{pa 4130 3260}%
\special{pa 4136 3266}%
\special{pa 4140 3270}%
\special{pa 4146 3276}%
\special{pa 4150 3280}%
\special{pa 4156 3286}%
\special{pa 4160 3290}%
\special{pa 4166 3296}%
\special{pa 4170 3300}%
\special{pa 4176 3306}%
\special{pa 4180 3310}%
\special{pa 4186 3316}%
\special{pa 4190 3320}%
\special{pa 4196 3326}%
\special{pa 4200 3330}%
\special{pa 4206 3336}%
\special{pa 4210 3340}%
\special{pa 4216 3346}%
\special{pa 4220 3350}%
\special{pa 4226 3356}%
\special{pa 4230 3360}%
\special{pa 4236 3366}%
\special{pa 4240 3370}%
\special{pa 4246 3376}%
\special{pa 4250 3380}%
\special{pa 4256 3386}%
\special{pa 4260 3390}%
\special{pa 4266 3396}%
\special{pa 4270 3400}%
\special{pa 4276 3406}%
\special{pa 4280 3410}%
\special{pa 4286 3416}%
\special{pa 4290 3420}%
\special{pa 4296 3426}%
\special{pa 4300 3430}%
\special{pa 4306 3436}%
\special{pa 4310 3440}%
\special{pa 4316 3446}%
\special{pa 4320 3450}%
\special{pa 4326 3456}%
\special{pa 4330 3460}%
\special{pa 4336 3466}%
\special{pa 4340 3470}%
\special{pa 4346 3476}%
\special{pa 4350 3480}%
\special{pa 4356 3486}%
\special{pa 4360 3490}%
\special{pa 4366 3496}%
\special{pa 4370 3500}%
\special{pa 4376 3506}%
\special{pa 4380 3510}%
\special{pa 4386 3516}%
\special{pa 4390 3520}%
\special{pa 4396 3526}%
\special{pa 4400 3530}%
\special{pa 4406 3536}%
\special{pa 4410 3540}%
\special{pa 4416 3546}%
\special{pa 4420 3550}%
\special{pa 4426 3556}%
\special{pa 4430 3560}%
\special{pa 4436 3566}%
\special{pa 4440 3570}%
\special{pa 4446 3576}%
\special{pa 4450 3580}%
\special{pa 4456 3586}%
\special{pa 4460 3590}%
\special{pa 4466 3596}%
\special{pa 4470 3600}%
\special{pa 4476 3606}%
\special{pa 4480 3610}%
\special{pa 4486 3616}%
\special{pa 4490 3620}%
\special{pa 4496 3626}%
\special{pa 4500 3630}%
\special{pa 4506 3636}%
\special{pa 4510 3640}%
\special{pa 4516 3646}%
\special{pa 4520 3650}%
\special{pa 4526 3656}%
\special{pa 4530 3660}%
\special{pa 4536 3666}%
\special{pa 4540 3670}%
\special{pa 4546 3676}%
\special{pa 4550 3680}%
\special{pa 4556 3686}%
\special{pa 4560 3690}%
\special{pa 4566 3696}%
\special{pa 4570 3700}%
\special{pa 4576 3706}%
\special{pa 4580 3710}%
\special{pa 4586 3716}%
\special{pa 4590 3720}%
\special{pa 4596 3726}%
\special{pa 4600 3730}%
\special{pa 4606 3736}%
\special{pa 4610 3740}%
\special{pa 4616 3746}%
\special{pa 4620 3750}%
\special{pa 4626 3756}%
\special{pa 4630 3760}%
\special{pa 4630 3760}%
\special{sp -0.045}%
% VECTOR 3 0 3 0
% 2 1680 2680 4810 1700
% 
\special{pn 4}%
\special{pa 1680 2680}%
\special{pa 4810 1700}%
\special{fp}%
\special{sh 1}%
\special{pa 4810 1700}%
\special{pa 4740 1702}%
\special{pa 4760 1716}%
\special{pa 4752 1740}%
\special{pa 4810 1700}%
\special{fp}%
% ELLIPSE 2 0 3 0
% 4 3100 1610 3950 1770 3950 1770 3950 1770
% 
%%%%%H^2_1の上の円  {
\special{pn 8}%
\special{ar 3100 1610 850 160  0.0000000 6.2831853}%
%%%%%%  }
% ELLIPSE 2 0 3 0
% 4 3100 2840 3940 3020 3940 3020 3940 3020
% 
\special{pn 8}%
\special{ar 3100 2840 840 180  0.0000000 6.2831853}%
% ELLIPSE 2 2 3 0
% 4 3100 1180 4120 1320 4120 1320 4120 1320
% 
\special{pn 8}%
\special{ar 3100 1180 1020 140  0.0000000 0.0206897}%
\special{ar 3100 1180 1020 140  0.0827586 0.1034483}%
\special{ar 3100 1180 1020 140  0.1655172 0.1862069}%
\special{ar 3100 1180 1020 140  0.2482759 0.2689655}%
\special{ar 3100 1180 1020 140  0.3310345 0.3517241}%
\special{ar 3100 1180 1020 140  0.4137931 0.4344828}%
\special{ar 3100 1180 1020 140  0.4965517 0.5172414}%
\special{ar 3100 1180 1020 140  0.5793103 0.6000000}%
\special{ar 3100 1180 1020 140  0.6620690 0.6827586}%
\special{ar 3100 1180 1020 140  0.7448276 0.7655172}%
\special{ar 3100 1180 1020 140  0.8275862 0.8482759}%
\special{ar 3100 1180 1020 140  0.9103448 0.9310345}%
\special{ar 3100 1180 1020 140  0.9931034 1.0137931}%
\special{ar 3100 1180 1020 140  1.0758621 1.0965517}%
\special{ar 3100 1180 1020 140  1.1586207 1.1793103}%
\special{ar 3100 1180 1020 140  1.2413793 1.2620690}%
\special{ar 3100 1180 1020 140  1.3241379 1.3448276}%
\special{ar 3100 1180 1020 140  1.4068966 1.4275862}%
\special{ar 3100 1180 1020 140  1.4896552 1.5103448}%
\special{ar 3100 1180 1020 140  1.5724138 1.5931034}%
\special{ar 3100 1180 1020 140  1.6551724 1.6758621}%
\special{ar 3100 1180 1020 140  1.7379310 1.7586207}%
\special{ar 3100 1180 1020 140  1.8206897 1.8413793}%
\special{ar 3100 1180 1020 140  1.9034483 1.9241379}%
\special{ar 3100 1180 1020 140  1.9862069 2.0068966}%
\special{ar 3100 1180 1020 140  2.0689655 2.0896552}%
\special{ar 3100 1180 1020 140  2.1517241 2.1724138}%
\special{ar 3100 1180 1020 140  2.2344828 2.2551724}%
\special{ar 3100 1180 1020 140  2.3172414 2.3379310}%
\special{ar 3100 1180 1020 140  2.4000000 2.4206897}%
\special{ar 3100 1180 1020 140  2.4827586 2.5034483}%
\special{ar 3100 1180 1020 140  2.5655172 2.5862069}%
\special{ar 3100 1180 1020 140  2.6482759 2.6689655}%
\special{ar 3100 1180 1020 140  2.7310345 2.7517241}%
\special{ar 3100 1180 1020 140  2.8137931 2.8344828}%
\special{ar 3100 1180 1020 140  2.8965517 2.9172414}%
\special{ar 3100 1180 1020 140  2.9793103 3.0000000}%
\special{ar 3100 1180 1020 140  3.0620690 3.0827586}%
\special{ar 3100 1180 1020 140  3.1448276 3.1655172}%
\special{ar 3100 1180 1020 140  3.2275862 3.2482759}%
\special{ar 3100 1180 1020 140  3.3103448 3.3310345}%
\special{ar 3100 1180 1020 140  3.3931034 3.4137931}%
\special{ar 3100 1180 1020 140  3.4758621 3.4965517}%
\special{ar 3100 1180 1020 140  3.5586207 3.5793103}%
\special{ar 3100 1180 1020 140  3.6413793 3.6620690}%
\special{ar 3100 1180 1020 140  3.7241379 3.7448276}%
\special{ar 3100 1180 1020 140  3.8068966 3.8275862}%
\special{ar 3100 1180 1020 140  3.8896552 3.9103448}%
\special{ar 3100 1180 1020 140  3.9724138 3.9931034}%
\special{ar 3100 1180 1020 140  4.0551724 4.0758621}%
\special{ar 3100 1180 1020 140  4.1379310 4.1586207}%
\special{ar 3100 1180 1020 140  4.2206897 4.2413793}%
\special{ar 3100 1180 1020 140  4.3034483 4.3241379}%
\special{ar 3100 1180 1020 140  4.3862069 4.4068966}%
\special{ar 3100 1180 1020 140  4.4689655 4.4896552}%
\special{ar 3100 1180 1020 140  4.5517241 4.5724138}%
\special{ar 3100 1180 1020 140  4.6344828 4.6551724}%
\special{ar 3100 1180 1020 140  4.7172414 4.7379310}%
\special{ar 3100 1180 1020 140  4.8000000 4.8206897}%
\special{ar 3100 1180 1020 140  4.8827586 4.9034483}%
\special{ar 3100 1180 1020 140  4.9655172 4.9862069}%
\special{ar 3100 1180 1020 140  5.0482759 5.0689655}%
\special{ar 3100 1180 1020 140  5.1310345 5.1517241}%
\special{ar 3100 1180 1020 140  5.2137931 5.2344828}%
\special{ar 3100 1180 1020 140  5.2965517 5.3172414}%
\special{ar 3100 1180 1020 140  5.3793103 5.4000000}%
\special{ar 3100 1180 1020 140  5.4620690 5.4827586}%
\special{ar 3100 1180 1020 140  5.5448276 5.5655172}%
\special{ar 3100 1180 1020 140  5.6275862 5.6482759}%
\special{ar 3100 1180 1020 140  5.7103448 5.7310345}%
\special{ar 3100 1180 1020 140  5.7931034 5.8137931}%
\special{ar 3100 1180 1020 140  5.8758621 5.8965517}%
\special{ar 3100 1180 1020 140  5.9586207 5.9793103}%
\special{ar 3100 1180 1020 140  6.0413793 6.0620690}%
\special{ar 3100 1180 1020 140  6.1241379 6.1448276}%
\special{ar 3100 1180 1020 140  6.2068966 6.2275862}%
% ELLIPSE 2 2 3 0
% 4 3100 3300 4150 3460 4150 3460 4150 3460
% 
% 漸近線の下の円{
\special{pn 8}%
\special{ar 3100 3300 1050 160  0.0000000 0.0198347}%
\special{ar 3100 3300 1050 160  0.0793388 0.0991736}%
\special{ar 3100 3300 1050 160  0.1586777 0.1785124}%
\special{ar 3100 3300 1050 160  0.2380165 0.2578512}%
\special{ar 3100 3300 1050 160  0.3173554 0.3371901}%
\special{ar 3100 3300 1050 160  0.3966942 0.4165289}%
\special{ar 3100 3300 1050 160  0.4760331 0.4958678}%
\special{ar 3100 3300 1050 160  0.5553719 0.5752066}%
\special{ar 3100 3300 1050 160  0.6347107 0.6545455}%
\special{ar 3100 3300 1050 160  0.7140496 0.7338843}%
\special{ar 3100 3300 1050 160  0.7933884 0.8132231}%
\special{ar 3100 3300 1050 160  0.8727273 0.8925620}%
\special{ar 3100 3300 1050 160  0.9520661 0.9719008}%
\special{ar 3100 3300 1050 160  1.0314050 1.0512397}%
\special{ar 3100 3300 1050 160  1.1107438 1.1305785}%
\special{ar 3100 3300 1050 160  1.1900826 1.2099174}%
\special{ar 3100 3300 1050 160  1.2694215 1.2892562}%
\special{ar 3100 3300 1050 160  1.3487603 1.3685950}%
\special{ar 3100 3300 1050 160  1.4280992 1.4479339}%
\special{ar 3100 3300 1050 160  1.5074380 1.5272727}%
\special{ar 3100 3300 1050 160  1.5867769 1.6066116}%
\special{ar 3100 3300 1050 160  1.6661157 1.6859504}%
\special{ar 3100 3300 1050 160  1.7454545 1.7652893}%
\special{ar 3100 3300 1050 160  1.8247934 1.8446281}%
\special{ar 3100 3300 1050 160  1.9041322 1.9239669}%
\special{ar 3100 3300 1050 160  1.9834711 2.0033058}%
\special{ar 3100 3300 1050 160  2.0628099 2.0826446}%
\special{ar 3100 3300 1050 160  2.1421488 2.1619835}%
\special{ar 3100 3300 1050 160  2.2214876 2.2413223}%
\special{ar 3100 3300 1050 160  2.3008264 2.3206612}%
\special{ar 3100 3300 1050 160  2.3801653 2.4000000}%
\special{ar 3100 3300 1050 160  2.4595041 2.4793388}%
\special{ar 3100 3300 1050 160  2.5388430 2.5586777}%
\special{ar 3100 3300 1050 160  2.6181818 2.6380165}%
\special{ar 3100 3300 1050 160  2.6975207 2.7173554}%
\special{ar 3100 3300 1050 160  2.7768595 2.7966942}%
\special{ar 3100 3300 1050 160  2.8561983 2.8760331}%
\special{ar 3100 3300 1050 160  2.9355372 2.9553719}%
\special{ar 3100 3300 1050 160  3.0148760 3.0347107}%
\special{ar 3100 3300 1050 160  3.0942149 3.1140496}%
\special{ar 3100 3300 1050 160  3.1735537 3.1933884}%
\special{ar 3100 3300 1050 160  3.2528926 3.2727273}%
\special{ar 3100 3300 1050 160  3.3322314 3.3520661}%
\special{ar 3100 3300 1050 160  3.4115702 3.4314050}%
\special{ar 3100 3300 1050 160  3.4909091 3.5107438}%
\special{ar 3100 3300 1050 160  3.5702479 3.5900826}%
\special{ar 3100 3300 1050 160  3.6495868 3.6694215}%
\special{ar 3100 3300 1050 160  3.7289256 3.7487603}%
\special{ar 3100 3300 1050 160  3.8082645 3.8280992}%
\special{ar 3100 3300 1050 160  3.8876033 3.9074380}%
\special{ar 3100 3300 1050 160  3.9669421 3.9867769}%
\special{ar 3100 3300 1050 160  4.0462810 4.0661157}%
\special{ar 3100 3300 1050 160  4.1256198 4.1454545}%
\special{ar 3100 3300 1050 160  4.2049587 4.2247934}%
\special{ar 3100 3300 1050 160  4.2842975 4.3041322}%
\special{ar 3100 3300 1050 160  4.3636364 4.3834711}%
\special{ar 3100 3300 1050 160  4.4429752 4.4628099}%
\special{ar 3100 3300 1050 160  4.5223140 4.5421488}%
\special{ar 3100 3300 1050 160  4.6016529 4.6214876}%
\special{ar 3100 3300 1050 160  4.6809917 4.7008264}%
\special{ar 3100 3300 1050 160  4.7603306 4.7801653}%
\special{ar 3100 3300 1050 160  4.8396694 4.8595041}%
\special{ar 3100 3300 1050 160  4.9190083 4.9388430}%
\special{ar 3100 3300 1050 160  4.9983471 5.0181818}%
\special{ar 3100 3300 1050 160  5.0776860 5.0975207}%
\special{ar 3100 3300 1050 160  5.1570248 5.1768595}%
\special{ar 3100 3300 1050 160  5.2363636 5.2561983}%
\special{ar 3100 3300 1050 160  5.3157025 5.3355372}%
\special{ar 3100 3300 1050 160  5.3950413 5.4148760}%
\special{ar 3100 3300 1050 160  5.4743802 5.4942149}%
\special{ar 3100 3300 1050 160  5.5537190 5.5735537}%
\special{ar 3100 3300 1050 160  5.6330579 5.6528926}%
\special{ar 3100 3300 1050 160  5.7123967 5.7322314}%
\special{ar 3100 3300 1050 160  5.7917355 5.8115702}%
\special{ar 3100 3300 1050 160  5.8710744 5.8909091}%
\special{ar 3100 3300 1050 160  5.9504132 5.9702479}%
\special{ar 3100 3300 1050 160  6.0297521 6.0495868}%
\special{ar 3100 3300 1050 160  6.1090909 6.1289256}%
\special{ar 3100 3300 1050 160  6.1884298 6.2082645}%
\special{ar 3100 3300 1050 160  6.2677686 6.2832853}%
%%%%% }
% STR 2 0 3 0
% 3 3750 2290 3750 2390 2 0
% $1$
\put(37.5000,-23.9000){\makebox(0,0)[lb]{$1$}}%
% STR 2 0 3 0
% 3 2210 2270 2210 2370 2 0
% $-1$
\put(22.1000,-23.7000){\makebox(0,0)[lb]{$-1$}}%
% STR 2 0 3 0
% 3 1440 1270 1440 1370 2 0
% {\LARGE $\mathbb{H}^2_1$}
\put(14.4000,-13.7000){\makebox(0,0)[lb]{{\LARGE $\mathbb{H}^2_1$}}}%
% LINE 2 0 3 0
% 2 890 300 890 610
% 
\special{pn 8}%
\special{pa 890 300}%
\special{pa 890 610}%
\special{fp}%
% LINE 2 0 3 0
% 2 890 610 580 610
% 
\special{pn 8}%
\special{pa 890 610}%
\special{pa 580 610}%
\special{fp}%
% STR 2 0 3 0
% 3 560 421 560 580 2 0
% {\LARGE $\mathbb{E}^3_2$}
\put(5.6000,-5.8000){\makebox(0,0)[lb]{{\LARGE $\mathbb{E}^3_2$}}}%
% STR 2 0 3 0
% 3 2070 790 2070 890 2 0
%
\pause
%%%%% x(M)の表の線
\special{pn 13}%
\special{pa 1540 780}%
\special{pa 1572 792}%
\special{pa 1602 806}%
\special{pa 1632 818}%
\special{pa 1664 830}%
\special{pa 1694 842}%
\special{pa 1724 852}%
\special{pa 1756 864}%
\special{pa 1786 876}%
\special{pa 1818 886}%
\special{pa 1848 896}%
\special{pa 1880 906}%
\special{pa 1910 916}%
\special{pa 1942 924}%
\special{pa 1972 932}%
\special{pa 2004 940}%
\special{pa 2034 948}%
\special{pa 2066 954}%
\special{pa 2098 958}%
\special{pa 2130 964}%
\special{pa 2162 968}%
\special{pa 2194 970}%
\special{pa 2226 972}%
\special{pa 2258 974}%
\special{pa 2290 974}%
\special{pa 2322 972}%
\special{pa 2354 970}%
\special{pa 2388 966}%
\special{pa 2420 962}%
\special{pa 2452 956}%
\special{pa 2486 950}%
\special{pa 2520 940}%
\special{pa 2552 930}%
\special{pa 2586 920}%
\special{pa 2620 908}%
\special{pa 2654 896}%
\special{pa 2688 884}%
\special{pa 2720 870}%
\special{pa 2754 858}%
\special{pa 2786 846}%
\special{pa 2820 834}%
\special{pa 2852 824}%
\special{pa 2884 814}%
\special{pa 2914 806}%
\special{pa 2944 800}%
\special{pa 2974 796}%
\special{pa 3004 794}%
\special{pa 3032 794}%
\special{pa 3058 798}%
\special{pa 3086 804}%
\special{pa 3110 814}%
\special{pa 3134 828}%
\special{pa 3158 844}%
\special{pa 3180 862}%
\special{pa 3204 884}%
\special{pa 3224 906}%
\special{pa 3246 930}%
\special{pa 3268 956}%
\special{pa 3290 982}%
\special{pa 3312 1008}%
\special{pa 3334 1034}%
\special{pa 3358 1060}%
\special{pa 3380 1086}%
\special{pa 3406 1110}%
\special{pa 3432 1132}%
\special{pa 3458 1152}%
\special{pa 3486 1170}%
\special{pa 3516 1186}%
\special{pa 3546 1200}%
\special{pa 3576 1214}%
\special{pa 3608 1224}%
\special{pa 3640 1234}%
\special{pa 3674 1242}%
\special{pa 3708 1248}%
\special{pa 3742 1252}%
\special{pa 3776 1254}%
\special{pa 3810 1256}%
\special{pa 3846 1256}%
\special{pa 3882 1254}%
\special{pa 3916 1252}%
\special{pa 3952 1246}%
\special{pa 3988 1240}%
\special{pa 4022 1234}%
\special{pa 4058 1226}%
\special{pa 4092 1216}%
\special{pa 4126 1204}%
\special{pa 4160 1192}%
\special{pa 4194 1180}%
\special{pa 4226 1166}%
\special{pa 4258 1150}%
\special{pa 4288 1134}%
\special{pa 4318 1116}%
\special{pa 4348 1098}%
\special{pa 4376 1080}%
\special{pa 4404 1060}%
\special{pa 4430 1038}%
\special{pa 4454 1016}%
\special{pa 4476 994}%
\special{pa 4480 990}%
\special{sp}%
% SPLINE 1 2 3 0
% 8 1560 780 1660 660 2540 340 3530 460 4200 510 4480 760 4500 960 4500 960
% 
\special{pn 13}%
\special{pa 1560 780}%
\special{pa 1580 756}%
\special{pa 1600 730}%
\special{pa 1620 706}%
\special{pa 1642 682}%
\special{pa 1664 658}%
\special{pa 1686 636}%
\special{pa 1710 614}%
\special{pa 1734 592}%
\special{pa 1758 572}%
\special{pa 1784 554}%
\special{pa 1810 536}%
\special{pa 1838 518}%
\special{pa 1866 502}%
\special{pa 1894 486}%
\special{pa 1924 472}%
\special{pa 1952 458}%
\special{pa 1982 444}%
\special{pa 2014 432}%
\special{pa 2044 422}%
\special{pa 2076 410}%
\special{pa 2108 400}%
\special{pa 2140 392}%
\special{pa 2172 384}%
\special{pa 2204 376}%
\special{pa 2238 368}%
\special{pa 2270 362}%
\special{pa 2304 358}%
\special{pa 2336 352}%
\special{pa 2370 350}%
\special{pa 2402 346}%
\special{pa 2436 344}%
\special{pa 2470 342}%
\special{pa 2502 340}%
\special{pa 2536 340}%
\special{pa 2568 340}%
\special{pa 2602 342}%
\special{pa 2634 344}%
\special{pa 2666 346}%
\special{pa 2698 348}%
\special{pa 2732 352}%
\special{pa 2764 354}%
\special{pa 2796 358}%
\special{pa 2828 364}%
\special{pa 2858 368}%
\special{pa 2890 372}%
\special{pa 2922 378}%
\special{pa 2954 382}%
\special{pa 2986 388}%
\special{pa 3018 394}%
\special{pa 3048 400}%
\special{pa 3080 404}%
\special{pa 3112 410}%
\special{pa 3142 416}%
\special{pa 3174 422}%
\special{pa 3206 426}%
\special{pa 3236 432}%
\special{pa 3268 436}%
\special{pa 3300 440}%
\special{pa 3332 444}%
\special{pa 3362 448}%
\special{pa 3394 452}%
\special{pa 3426 454}%
\special{pa 3458 458}%
\special{pa 3488 458}%
\special{pa 3520 460}%
\special{pa 3552 460}%
\special{pa 3584 462}%
\special{pa 3616 460}%
\special{pa 3648 460}%
\special{pa 3680 460}%
\special{pa 3712 458}%
\special{pa 3746 458}%
\special{pa 3778 456}%
\special{pa 3810 456}%
\special{pa 3842 456}%
\special{pa 3874 456}%
\special{pa 3906 456}%
\special{pa 3938 458}%
\special{pa 3970 460}%
\special{pa 4002 462}%
\special{pa 4034 468}%
\special{pa 4066 472}%
\special{pa 4096 478}%
\special{pa 4128 486}%
\special{pa 4160 496}%
\special{pa 4192 506}%
\special{pa 4222 520}%
\special{pa 4252 534}%
\special{pa 4282 550}%
\special{pa 4312 568}%
\special{pa 4340 586}%
\special{pa 4366 606}%
\special{pa 4392 628}%
\special{pa 4414 652}%
\special{pa 4434 678}%
\special{pa 4452 704}%
\special{pa 4468 730}%
\special{pa 4480 760}%
\special{pa 4488 790}%
\special{pa 4494 820}%
\special{pa 4498 852}%
\special{pa 4500 884}%
\special{pa 4500 918}%
\special{pa 4500 950}%
\special{pa 4500 960}%
\special{sp -0.045}%
%
% ${\bf x}(M)$
\put(20.7000,-8.9000){\makebox(0,0)[lb]{$M$}}%
% VECTOR 1 0 3 0
% 2 3100 2230 3550 1230
% 
% STR 2 0 3 0
% 3 3500 1090 3500 1190 2 0
% ${\bf x}(p)$
\put(35.0000,-11.9000){\makebox(0,0)[lb]{${\bf x}(p)$}}%
% DOT 0 0 3 0
% 2 3550 1200 3550 1200
% 
\special{pn 20}%
\special{sh 1}%
\special{ar 3550 1200 10 10 0  6.28318530717959E+0000}%
\special{sh 1}%
\special{ar 3550 1200 10 10 0  6.28318530717959E+0000}%
\pause
%%%  接ベクトルと位置ベクトル {
\special{pn 13}%
\special{pa 3100 2230}%
\special{pa 3550 1230}%
\special{fp}%
\special{sh 1}%
\special{pa 3550 1230}%
\special{pa 3504 1284}%
\special{pa 3528 1280}%
\special{pa 3542 1300}%
\special{pa 3550 1230}%
\special{fp}%
% VECTOR 1 0 3 0
% 2 3520 1190 2980 960
% 
\special{pn 13}%
\special{pa 3520 1190}%
\special{pa 2980 960}%
\special{fp}%
\special{sh 1}%
\special{pa 2980 960}%
\special{pa 3034 1006}%
\special{pa 3030 982}%
\special{pa 3050 968}%
\special{pa 2980 960}%
\special{fp}%
%}
% STR 2 0 3 0
% 3 2680 1160 2680 1260 2 0
% ${\bf x}_*(e_1^p)$
\put(26.8000,-12.6000){\makebox(0,0)[lb]{$e_1^p$}}%
% POLYGON 2 0 3 0
% 6 2880 460 4144 913 3220 2980 1834 2527 1834 2527 2880 460
% 
\pause
\color{red}
\special{pn 8}%
\special{pa 2880 460}%
\special{pa 4144 914}%
\special{pa 3220 2980}%
\special{pa 1834 2528}%
\special{pa 1834 2528}%
\special{pa 2880 460}%
\special{fp}%
% STR 2 0 3 0
% 3 550 1840 550 1940 2 0
% $\tilde{\nu}(p)={\bf x}(p)\wedge {\bf x}_*(e_1^p)$
\put(5.5000,-19.4000){\makebox(0,0)[lb]{$\tilde{\nu}(p)={\bf x}(p)\wedge e_1^p$}}%
% SPLINE 2 0 3 0
% 9 1850 1880 2026 1949 2043 2010 1975 2014 1963 1986 2059 1922 2122 1935 2224 1990 2224 1990
% 
%%% くるん曲線 {
\special{pn 8}%
\special{pa 1850 1880}%
\special{pa 1884 1884}%
\special{pa 1916 1888}%
\special{pa 1948 1894}%
\special{pa 1978 1908}%
\special{pa 2004 1926}%
\special{pa 2028 1952}%
\special{pa 2046 1986}%
\special{pa 2042 2012}%
\special{pa 2010 2024}%
\special{pa 1976 2014}%
\special{pa 1964 1986}%
\special{pa 1976 1958}%
\special{pa 2002 1936}%
\special{pa 2038 1926}%
\special{pa 2074 1922}%
\special{pa 2106 1930}%
\special{pa 2134 1940}%
\special{pa 2164 1954}%
\special{pa 2192 1970}%
\special{pa 2218 1988}%
\special{pa 2224 1990}%
\special{sp}%
%%%%%%%}
% STR 2 0 3 0
% 3 2140 2080 2140 2180 2 0
% {\LARGE $\circlearrowleft$}
\put(21.4000,-21.8000){\makebox(0,0)[lb]{{\LARGE $\circlearrowleft$}}}%
% SPLINE 1 0 3 0
% 6 1540 780 2520 940 3100 810 3430 1130 4470 1000 4480 990
% 
\end{picture}%
}
% \end{center}
% }


\subsection{in \(\mathbb{S}^3_1\) or \(\mathbb{H}^3_1\)}

\frame{
  \frametitle{\insertsection \quad \quad{\normalsize $\cdots$ \insertsubsection}}
%\begin{center}
\vspace{-5pt}
\hspace{0pt}
%\begin{minipage}[c]{0.9\textwidth}
%\scalebox{0.9}{
\begin{thm}[D. S. Kim--Y. H. Kim]
$M^2_1\subset \mathbb{S}^3_1$ (または$\mathbb{H}^3_1)$ : 向きづけられたローレンツ超曲面
\begin{center}
\textcolor<2->{red}{形作用素が対角化不可能 かつ\, $H$,$K$ が一定}\vspace{3pt}\\
\rotatebox{90}{{\large $\Leftrightarrow$}} \vspace{-2pt}\\
$M^2_1$はB-scrollまたはcomplex circleどちらかの開部分 \vspace{-0pt}
\end{center}
\end{thm}
%}
%\end{minipage}
%\end{center}
\pause
\begin{center}
\raisebox{-2mm}[0mm]{\includegraphics[width=15pt, height=20pt, scale=0.2]{yazirusitate.png}} 
\vspace{-8pt}
\end{center}
\begin{table}
\scalebox{0.95}{
\begin{tabular}{|l||l|c|c|c|c|} \hline 
& \multicolumn{2}{c|}{\raisebox{-1mm}{\shortstack{\rule{0pt}{0.5pt}\\ $M^2_1$\\ \rule{0pt}{0.5pt}}}} & \raisebox{1mm}{$\tilde{\nu}$} & \raisebox{1mm}{$K$} & \raisebox{1mm}{$H$} \vspace{0pt} \\ \hline \hline %\cline{2-6} %\cline{1-6} 
& \uncover<3->{\raisebox{-0mm}{\shortstack{\rule{0pt}{1pt}\\ \textcolor<3->{red}{(Main Result 1)} \,\, $S^1_{\mathbb{C}}(\kappa)$}}} & \uncover<3->{\raisebox{-0mm}{$\kappa =-1$}} & \uncover<3->{\raisebox{-0mm}{$1$-type}} & \uncover<3->{\raisebox{-0mm}{$0$}} & \uncover<3->{\raisebox{-0mm}{$0$}} \\ \cline{3-6}
\uncover<3->{in \,$\mathbb{H}^3_1$} & \uncover<3->{\hspace{15pt}complex circle} & \uncover<3->{\raisebox{0mm}{$\kappa \neq -1$}} & \uncover<3->{$\infty$-type} & \uncover<3->{\raisebox{-0mm}{$0$}} & \uncover<3->{\raisebox{-0mm}{$\neq0$}} \\ \cline{2-6}
& \uncover<3->{\raisebox{-0mm}{\shortstack{\rule{0pt}{1pt}\\ \textcolor<3->{red}{(Main Result 2)} \,\, $\mathcal{B}(k_2)$}}} & \uncover<3->{$k_2=\pm1$} & \uncover<3->{\raisebox{-0mm}{$\infty$-type}} & \uncover<3->{\raisebox{-0mm}{$0$}} & \uncover<3->{\raisebox{-0mm}{$\neq0$}} \\ \cline{3-6}
& \uncover<3->{\raisebox{-0mm}{\hspace{30pt}B-scroll}} & \uncover<3->{$k_2 \neq \pm1$} & \uncover<3->{null $2$-type} & \uncover<3->{\raisebox{-0mm}{$\neq0$}} & \uncover<3->{\raisebox{-0mm}{$\neq0$}} \\ \cline{1-6} 
\raisebox{2mm}{in \,$\mathbb{S}^3_1$}\vspace{0pt} & \raisebox{-0mm}{\shortstack{\rule{0pt}{1pt}\\ (B--C--D 2017) \hspace{7pt} $\mathcal{B}(k_2)$\\ \hspace{-8pt}B-scroll}}  & & \raisebox{2.5mm}{null $2$-type}\vspace{-0pt} & \raisebox{2.5mm}{$\neq 0$} & \raisebox{2.5mm}{$\neq0$} \\ \hline
%& \raisebox{0mm}{\textcolor<7->{red}{B-scroll}} & & & & \\ \hline
\end{tabular}
}
\end{table}
}




\section{Part 1-2 \quad B-scroll}

\subsection{定義}

\frame{
  \frametitle{\insertsection \quad \quad{\normalsize $\cdots$ \insertsubsection}}

    \hspace{18pt}\(\gamma\) \hspace{14pt}: \(\mathbb{S}^{3}_1\)または\(\mathbb{H}^{3}_1\)のnull曲線

    \((A, B, C)\) : \(\gamma\)上の\textcolor{red}{Cartan frame field}
    \begin{center}
    i.e. \,
    \(
    \left\{
    \begin{array}{l}
    \langle A,A\rangle = \langle B,B\rangle=0,\quad \langle A,B\rangle=-1, \\
    \langle A,C\rangle = \langle B,C\rangle=0,\quad \langle C,C\rangle=1,\\
    \dot{\gamma}(s)=A(s), \\
    \dot{A}(s)=k_1(s)C(s), \\
    \dot{C}(s)=\textcolor<3->{red}{k_2(s)} A(s)+k_1(s)B(s),\\
    \dot{B}(s)=\textcolor<3->{red}{k_2(s)} C(s)+\varepsilon \gamma(s). 
    \end{array}
    \right.
    \)
    \end{center}
    \vspace{-6pt}
    \pause
  \begin{defi}
    \(M\)を，次のようにパラメータづけされたローレンツ曲面とする:
    \vspace{-7pt}
    \begin{table}
    \begin{tabular}{lll}
    \({\bf x} : M \hookrightarrow \mathbb{S}^{3}_1 \,\text{or}\, \mathbb{H}^{3}_1\)
    & \(\overset{\mathrm{def}}{\Leftrightarrow}\)
    & \({\bf x}(s, t) := \gamma(s)+tB(s)\)
    \vspace{-7pt}
    \end{tabular}
    \end{table}
    このとき，\vspace{-10pt}
    \begin{table}
      \begin{tabular}{lll}
      \(M\) : \textcolor{red}{{\it \(\gamma\)上のB-scroll}}
      & \(\overset{\mathrm{def}}{\Leftrightarrow}\)
      & \textcolor<3->{red}{\(k_2\)} : const
      \end{tabular}
      \end{table}
      \vspace{-8pt}
  \end{defi}
}



\frame{
  \frametitle{\insertsection \quad \quad{\normalsize $\cdots$ \insertsubsection}}
  \begin{table}
  \begin{tabular}{lr}
    \hspace{-50pt}
  \begin{minipage}{0.35\textwidth}
\scalebox{0.7}{\input{pic_B-scroll-10}}
  \end{minipage}
  \hspace{30pt}
&
\begin{minipage}{0.65\textwidth}
  \underline{Note}
\begin{itemize}
\item null曲線\(\gamma\)と\(\gamma\)に沿うFrenet型フレーム場から構成される線織面
\item \(2\)次元非退化ローレンツ曲面
\item 形作用素は対角化不可能，実固有値
\end{itemize}
\end{minipage}
\end{tabular}
\end{table}
}








\section{Part 2-1 \quad B-scrollの次元一般化}

% \frame{
%   \begin{center}
%     {\Large {\bf Part \Rnum{2}}}
%   \end{center}
% \begin{beamercolorbox}[
%   %wd=60mm, sep=2pt,
%   %shadow=true, 
%   rounded=true
%   ]{frametitle}
%   \begin{center}
%     {\Large generalizations of B-scroll in \(\tilde{M}^{m}_s\)}
%   \end{center}
%   \end{beamercolorbox}
% }

\subsection{定義}







\frame{
\frametitle{\insertsection \quad{\normalsize $\cdots$ \insertsubsection}}
\(\tilde{M}^{n+1}_1 =\) \textcolor{red}{\(\mathbb{E}^{n+1}_1\) or \(\mathbb{S}^{n+1}_1\) or \(\mathbb{H}^{n+1}_1\)}
\vspace{3pt}

\(M^n_1\) \hspace{9pt}: \(\tilde{M}^{n+1}_1\)のローレンツ超曲面
\vspace{3pt}

\hspace{3pt}\(A\) \hspace{14pt}: \(M^n_1\)の形作用素, \textcolor<4->{red}{対角化不可能, 実固有値}
\vspace{3pt}

\(P(x)\) \hspace{4pt}: \(A\)の最小多項式
%\vspace{10pt}
\begin{thm}[M. A. Magid]
\(A\)の固有値の個数が1 or 2個 かつ \(0\)でない固有値が1つ以下
\vspace{8pt}

\hspace{100pt}
\rotatebox{90}{{\Large \(\Leftarrow\)}}
\vspace{8pt}

\(M^n_1\)は次のいずれか：\vspace{8pt}
\begin{enumerate}
  \item \(M^n_1\) : generalized cylinder \uncover<2->{\quad \(\cdots\) ゼロ固有値を持つ}
  \item \(M^n_1\) : generalized umbilical hypersurface \uncover<2->{\quad \(\cdots\) ゼロ固有値を持たない}
\end{enumerate}
\vspace{-0pt}
\uncover<3->{
\hspace{60pt}
  \(
    \left\{
  \begin{aligned}
    & \, \text{degree} \, 2 \textcolor<4->{red}{\quad P(x)=(x-a)^2} \uncover<4->{\quad \textcolor{red}{\leftarrow \text{B-scrollはこの型}}}\\
    & \, \text{degree} \, 3 \quad P(x)=(x-a)^3
  \end{aligned}
    \right.
  \)
    }
\end{thm}
}


\subsection{形作用素}


\frame{
  \frametitle{\insertsection}
\(\tilde{M}\)が指数\(1\)のとき, 形作用素\(A\)はある基底に関して次の\(4\)つのいずれかの形をとる.
\[
  \begin{aligned}
  &\mbox{(I)}
  \,\,
  \begin{pmatrix}
  a_1 &       & 0 \\
      &\ddots &   \\
  0   &       & a_n    
  \end{pmatrix}
  ,
  &&
  \textcolor{red}{
  \mbox{(II)}
  \,\,
  \begin{pmatrix}
  a_0 & 0   & & & \\
  1   & a_0 & & & \hspace{-15pt}  \\
      &     & a_1 & & \\
      &  & & \ddots & \\
      &          & &        & a_{n-2}
  \end{pmatrix}
  ,
  }
  \\
  &\mbox{(III)}
  \,\,
  \begin{pmatrix}
  a_0 & 0   & 0   & & & \\
  0   & a_0 & 1   & & & \\
  -1  & 0   & a_0 & & & \\
      &     &     & a_1 & & \\
      &     &     &     & \ddots & \\
      &     &     &     &        & a_{n-3}
  \end{pmatrix}
  ,
  &&\mbox{(IV)}
  \,\,
  \begin{pmatrix}
  a_0 & b_0   & & & \\
  -b_0   & a_0 & & & \hspace{-15pt}  \\
      &     & a_1 & & \\
      &  & & \ddots & \\
      &          & &        & a_{n-2}
  \end{pmatrix}
  ,
  \end{aligned}
  \]
  ここで，$b_0\neq0$. 

}






\section{Part 2-2 \quad B-scrollの次元及び指数一般化}


\subsection{形作用素}


\frame{
  \frametitle{\insertsection \quad{\normalsize $\cdots$ \insertsubsection}}
\underline{\quad$\tilde{M}^m_2$の形作用素\quad} \\
\vspace{5pt}
{\footnotesize
\textcolor{red}{
$
\left(
{\renewcommand\arraystretch{1}
{\arraycolsep = 1mm
\begin{array}{cc|ccc}
\lambda & 0 & & & \\
 1 & \lambda &  & \\ \hline
 & &  \lambda &   & \\
 & &  & \ddots & \\
 & &    &   & \lambda 
\end{array}
}}
\right)
$
}
}
 ,\quad
{\footnotesize
$
\left(
{\renewcommand\arraystretch{0.5}
{\arraycolsep = 1mm
\begin{array}{cc|cc|ccc}
\lambda & 0 & & & & & \\
1 & \lambda & & & & & \\ \hline
 & & \lambda & 0 & & & \\
 & & 1 & \lambda & & & \\ \hline
 & &    &   & \lambda & & \\
 & &    &   &   & \ddots & \\
 & &    &   &   &        & \lambda
\end{array}
}}
\right)
$
}
 ,\quad
{\footnotesize
$
\left(
{\renewcommand\arraystretch{0.5}
{\arraycolsep = 1mm
\begin{array}{ccc|ccc}
\lambda &  &  &  & & \\
 1&  \lambda &  &  & & \\
 & 1& \lambda  &  & & \\ \hline
 &  &  & \lambda &  & \\
 &  &  &   & \ddots &  \\
 &  &  &   &   & \lambda  
\end{array}
}}
\right)
$
}
 ,\quad
\vspace{10pt} \\
{\footnotesize
$
\left(
{\renewcommand\arraystretch{0.8}
{\arraycolsep = 1mm
\begin{array}{ccc|ccc|ccc}
\lambda &         &         &   &   &   &   &        &   \\
1       & \lambda &         &   &   &   &   &        &   \\
        &  1      & \lambda &   &   &   &   &        &   \\ \hline
        &         &         & \lambda &         &  &   &        &   \\ 
        &         &         &   1     & \lambda &   &   &        &   \\
        &         &         &         &   1     & \lambda &   &        &   \\ \hline
 &   &   &   &   &   & \lambda &        &   \\
 &   &   &   &   &   &   & \ddots &   \\
 &   &   &   &   &   &   &        & \lambda \\
\end{array}
}}
\right)
$ ,\quad
{\footnotesize
$
\left(
{\renewcommand\arraystretch{0.5}
{\arraycolsep = 1mm
\begin{array}{cccc|ccc}
\lambda & 0 & & & & & \\
1 & \lambda & & & & & \\
 & 1 & \lambda & 0 & & & \\
 & & 1 & \lambda & & & \\ \hline
 & &    &   & \lambda & & \\
 & &    &   &   & \ddots & \\
 & &    &   &   &        & \lambda
\end{array}
}}
\right)
$
}
 ,\quad
$
\left(
{\renewcommand\arraystretch{0.5}
{\arraycolsep = 1mm
\begin{array}{ccccc|ccc}
\lambda&          &           &          &   &   &   &          \\
1      &  \lambda &           &          &   &          &          &   \\
       &      1   & \lambda   &          &   &   &   &          \\
       &          &  1        &  \lambda &   &  &   &          \\
       &          &           &    1     & \lambda &   &          &   \\ \hline
       &          &           &          &      & \lambda &        &   \\
       &          &           &          &   &     & \ddots &   \\
       &          &           &          &   &   &   &         \lambda  
\end{array}
}}
\right)
$
}
\pause
\begin{alertblock}{Main Result 3 (H. Kobayashi)}
  degree \(2\) のgeneralized umbilical hypersurface in \textcolor{red}{\(\mathbb{S}^{n+1}_2\) or \(\mathbb{H}^{n+1}_2\)} の具体例を構成
\end{alertblock}
}






% \frame{
% \frametitle{\insertsection \quad{\normalsize $\cdots$ \insertsubsection}}
% \(\tilde{M}^{n+1}_1\) : ローレンツ多様体，\(\mathrm{dim}\tilde{M}=n+1\)
% \vspace{3pt}

% \(M^n_1\) \hspace{9pt}: \(\tilde{M}^{n+1}_1\)のローレンツ超曲面，\(\mathrm{dim}M=n\)
% \vspace{3pt}

% \hspace{3pt}\(A\) \hspace{14pt}: \(M^n_1\)の形作用素, 対角化不可能
% \begin{defi}
%   \begin{table}
%     \begin{tabular}{lll}
%       \(M^n_1\) : \textcolor{red}{{\it generalized umbilical hypersurface}}
%       & \(\overset{\mathrm{def}}{\Leftrightarrow}\)
%       & \(A\)が\(0\)でない唯一の\\
%       &
%       &実固有値を持つ
%     \end{tabular}
%   \end{table}
% \(A\)の最小多項式が\(P(x)=(x-a)^2\)のとき，\(M^n_1\)を\textcolor{red}{degree \(2\)のgeneralized umbilical hypersurface}という (\(a\in \mathbb{R}\) : const)．
% \end{defi}

% \pause
% \begin{alertblock}{Main Result 3 (H. Kobayashi--N. Koike)}
%       {\it
%       \(M^n_1\) : degree \(2\)の g. u. h. , \quad \(S\) : \(M^n_1\)のスカラー曲率
%       \vspace{-3pt}
%       \begin{table}
%       \begin{tabular}{lll}
%        $S=0$ \hspace{12pt} & $\Rightarrow$ &  \hspace{-10pt} $\tilde{\nu}$ : infinite type \\
%        $S\neq 0$(一定) & $\Rightarrow$& \hspace{-10pt} $\tilde{\nu}$ : null $2$-type
%       \end{tabular}
%       \end{table}
%       \vspace{-10pt}
%       }
%       \end{alertblock}
%}




% \frame{
%   \frametitle{\insertsection \quad {\normalsize $\cdots$ \insertsubsection}}
%   \begin{minipage}{0.9\textwidth}
%     Cartan frame (E. Cartan) : \(\mathbb{E}^3_1\)におけるnull曲線に沿う\\
%     \hspace{122pt} Frenet型フレーム
%     \vspace{-5pt}
%     \begin{flushleft}
%       \hspace{60pt}
%       \raisebox{-2mm}[0mm]{\includegraphics[width=15pt, height=20pt, scale=0.2]{yazirusitate.png}} 
%       \hspace{10pt}
%       {\footnotesize 一般次元ローレンツ多様体に拡張}
%       \vspace{-5pt}
%       \end{flushleft}
%     general Frenet frame (K. L. Duggal--A. Bejancu)
%     \vspace{-5pt}
%     \uncover<2->{
%     \begin{flushleft}
%       \hspace{60pt}
%       \raisebox{-2mm}[0mm]{\includegraphics[width=15pt, height=20pt, scale=0.2]{yazirusitate.png}} 
%       \hspace{10pt}
%       {\footnotesize よりシンプルな形に再構成}
%       \vspace{-5pt}
%       \end{flushleft}
%       natural Frenet frame (D. H. Jin)
%       \vspace{-5pt}
%       \begin{flushleft}
%         \hspace{60pt}
%         \raisebox{-2mm}[0mm]{\includegraphics[width=15pt, height=20pt, scale=0.2]{yazirusitate.png}} 
%         \hspace{10pt}
%         {\footnotesize 特殊なパラメータをとることでさらにシンプルに}
%         \vspace{-5pt}
%         \end{flushleft}
%       Cartan frame (A. Ferr\'{a}ndez--A. Gim\'{e}nez--P. Lucas)
%     }
%   \end{minipage}
%   \begin{minipage}{0.09\textwidth}
%     \uncover<5->{
%    \raisebox{-48mm}[0mm]{\includegraphics[height=130pt, scale=0.5]{ko-yazirushi.png}} 
%     }
%   \end{minipage}
%   \vspace{-2pt}
%   \uncover<3->{
%   \begin{itemize}
%     \item これらは構成の際にscreen vercor bundleなどの概念が必要
%   \vspace{7pt}
%   \end{itemize}
%   }
%   \uncover<4->{
%   \underline{一方，H. Kobayashi--N. Koikeによる構成方法は}
%   \begin{itemize}
%     \item general Frenet frameを経由せず構成
%     \item 高度な概念必要なし
%     \item \(\gamma\)のパラメータに条件なし
%   \end{itemize}
%   }
% }










% \frame{
%   \frametitle{\insertsection \quad{\normalsize $\cdots$ \insertsubsection}}
%   \underline{\(\tilde{M}^n_1\)の形作用素} \vspace{5pt}

% \(\tilde{M}\)が指数\(1\)のとき, \(A_N\)はある基底に関して次の\(4\)つのいずれかの形を

% とる( $b_0\neq0$ ).
% \vspace{-5pt}
% \[
%   \begin{aligned}
%   &\mbox{(I)}
%   \,\,
%   \begin{pmatrix}
%   a_1 &       & 0 \\
%       &\ddots &   \\
%   0   &       & a_n    
%   \end{pmatrix}
%   ,
%   &&
%   \textcolor{red}{
%   \mbox{(II)}
%   \,\,
%   \begin{pmatrix}
%   a_0 & 0   & & & \\
%   1   & a_0 & & & \hspace{-15pt}  \\
%       &     & a_1 & & \\
%       &  & & \ddots & \\
%       &          & &        & a_{n-2}
%   \end{pmatrix}
%   ,}
%   \\
%   &\mbox{(III)}
%   \,\,
%   \begin{pmatrix}
%   a_0 & 0   & 0   & & & \\
%   0   & a_0 & 1   & & & \\
%   -1  & 0   & a_0 & & & \\
%       &     &     & a_1 & & \\
%       &     &     &     & \ddots & \\
%       &     &     &     &        & a_{n-3}
%   \end{pmatrix}
%   ,
%   &&\mbox{(IV)}
%   \,\,
%   \begin{pmatrix}
%   a_0 & b_0   & & & \\
%   -b_0   & a_0 & & & \hspace{-15pt}  \\
%       &     & a_1 & & \\
%       &  & & \ddots & \\
%       &          & &        & a_{n-2}
%   \end{pmatrix}
%   ,
%   \end{aligned}
%   \]
% }








\subsection{Frenet型フレーム場}





% \frame{
%   \frametitle{\insertsection \quad{\normalsize $\cdots$ \insertsubsection}}
% Cartan frame (E. Cartan) : \(\mathbb{E}^3_1\)におけるnull曲線に沿う\\
% \hspace{122pt} Frenet型フレーム
% \vspace{-3pt}
% \begin{flushleft}
%   \hspace{60pt}
%   \raisebox{-2mm}[0mm]{\includegraphics[width=15pt, height=18pt, scale=0.2]{yazirusitate.png}} 
%   \hspace{10pt}
%   {\small 一般次元ローレンツ多様体\(M^n_1\)に拡張}
%   \vspace{-2pt}
%   \end{flushleft}
% general Frenet frame (K. L. Duggal--A. Bejancu)
% \vspace{-3pt}
% \begin{flushleft}
%   \hspace{60pt}
%   \raisebox{-2mm}[0mm]{\includegraphics[width=15pt, height=18pt, scale=0.2]{yazirusitate.png}} 
%   \hspace{10pt}
%   {\small よりシンプルな形に再構成}
%   \vspace{-2pt}
%   \end{flushleft}
%   natural Frenet frame (D. H. Jin)
%   \pause
%   \vspace{-3pt}
%   \begin{flushleft}
%     \hspace{60pt}
%     \raisebox{-2mm}[0mm]{\includegraphics[width=15pt, height=18pt, scale=0.2]{yazirusitate.png}} 
%     \hspace{10pt}
%     {\small 指数\(2\)に拡張}
%     \vspace{-2pt}
%     \end{flushleft}
%   natural Frenet frame with index \(2\) (K. L. Duggal--A. Bejancu--D. H. Jin)
%   \vspace{15pt}
%    \pause

%   \underline{別証明}
%   \vspace{-5pt}
%   \[
%     \left\{
%   \begin{aligned}
%     & \nabla_{\dot{\gamma}}A \text{\,: non-null かつ \(^\forall Z_i\) : non-nullのとき，Cartan frame (H. Kobayashi)}\\
%     & \nabla_{\dot{\gamma}}A \text{\,: nullのとき，bi-null Cartan frame (M. Sakaki--A. U\c{c}um--K. \.{I}larslan)}
%   \end{aligned}
%     \right.
%   \]
%   \pause
%   \vspace{-5pt}
%   \begin{alertblock}{Main Result 4 (H. Kobayashi)}
%   degree \(2\) のgeneralized umbilical hypersurface in \(\mathbb{S}^{n+1}_2\) or \(\mathbb{H}^{n+1}_2\) の具体例を構成
%   \end{alertblock}
% }








\section{Part 2-3 \quad B-scrollの指数及び余次元一般化}



\subsection{in \(\mathbb{E}^m_1\)}

\frame{
  \frametitle{\insertsection \quad \quad{\normalsize $\cdots$ \insertsubsection}}
  \vspace{3pt}

\((A, B, C, Z_1, \ldots, Z_{m-3})\) : \(\gamma\)に沿うCartan frame field
\vspace{-3pt}
\begin{table}
  \begin{tabular}{cll}
    \(M\) : \textcolor{red}{generalized B-scroll}
    & \(\overset{\mathrm{def}}{\Leftrightarrow}\)
    & \(
      M \,:\, {\bf x}(s, t)=\gamma(s)+tB(s)  
      \)
  \end{tabular}
\end{table}
\vspace{0pt}
\pause

\underline{Note}
\begin{itemize}
\item null曲線\(\gamma\)と\(\gamma\)に沿うFrenet型フレーム場から構成される線織面
\item \(2\)次元非退化ローレンツ曲面
\item 形作用素は対角化不可能，実固有値
\end{itemize}
\pause

\begin{theorem}[D. S. Kim--Y. H. Kim--D. W. Yoon]
  % \begin{table}
  % \begin{tabular}{ll}
  % \begin{minipage}{0.42\textwidth}
    \vspace{0pt}

\hspace{3pt}\(M\)\hspace{2pt} : \textcolor{red}{\(\mathbb{E}^m_1\)}のnull scroll \vspace{3pt}

\hspace{3pt}\(H\)\hspace{2pt} : \(M\)の平均曲率ベクトル場\vspace{3pt}

\(A_H\) : \(M\)の\(H\)方向の形作用素
  % \end{minipage}
  % &
  % \uncover<3->{
    
  % \begin{minipage}{0.49\textwidth}
  %   \color<3, 6>{red}{
  %     {\small ※\(M\)は\(2\)次元非退化ローレンツ線織面}
  %   \begin{screen}
  %     \vspace{-20pt}
  %      \[M :\, {\bf x}(s, t)=\gamma(s)+tB(s) \vspace{-10pt}\]
  %     %  {\small i.e. \(M\)は\(\gamma\)と\(\gamma\)に沿うFrenet型フレーム場から構成される線織面}
  %      \vspace{-15pt}
  %   \end{screen}

  %   \hspace{70pt} 
  %   \rotatebox{90}{\(\Leftrightarrow\)}
  %   {\rm {\footnotesize \raisebox{1mm}{def}}}
  %   }
  % \end{minipage}
  % }
  % \end{tabular}
  % \vspace{-20pt}
  % \end{table}

\begin{table}
  \begin{tabular}{cll}
    \(A_H\)の最小多項式が\((x-a^2)^2\)
    & \hspace{3pt}\(\Leftrightarrow\)
    & \(M\) : generalized B-scroll\\
    (\(a\in \mathbb{R}\) : const)
    &
    &
  \end{tabular}
  % \uncover<4->{
%   \vspace{-15pt}
% \end{table}
% \begin{table}
%   \begin{tabular}{cll}
%     \(A_H\)の最小多項式が\((x-\textcolor<5>{red}{k_2^2})^2\)
%     & \(\Leftarrow\)
%     & \(M\) : \textcolor<6>{red}{generalized B-scroll}
%   \end{tabular}
\end{table}
  %}
\vspace{-5pt}
\end{theorem}
% \uncover<3->{
%   \underline{\(\gamma \subset \mathbb{E}^m_1\)に沿うCartan frame field}
%   \vspace{-5pt}
% \begin{center}
% \begin{minipage}{0.7\textwidth}
%   % \underline{\(\gamma\) in \(\mathbb{E}^m_1\)}
%   {\small
%   {\renewcommand\arraystretch{0.3}
%   {\arraycolsep = 0.5mm
%   %\hspace{-30pt}
%   \(
%   \begin{array}{c}
%   A \\ B\\ C\\ Z_1\\ \vdots \vspace{2pt}\\ Z_m
%   \end{array}
%   \)
%   }
%   {\renewcommand\arraystretch{0.6}
%   {\arraycolsep = 0.8mm
%   \(
%   \left(
%   \begin{array}{ccccc}
%   0 & -1&   &   &   \\
%   -1& 0 &   &   &   \\
%     &   & 1 &   &   \\
%     &   &   & \ddots &  \\
%     &   &   &   & 1 
%   \end{array}
%   \right)
%   \)
%   }}
%   } 
%   \, and\quad
%   %\vspace{-15pt}
%   \(
%   \left \{
%   \begin{array}{l}
%   \dot{A}=k_1C\\
%   \dot{C}=\textcolor<5>{red}{k_2}A+k_1B\\
%   \dot{B}=\textcolor<5>{red}{k_2} C+Z_1 \\
%   \dot{Z}_1=k_3A+k_4Z_2 \\
%   \dot{Z}_2=-k_4Z_1+k_5Z_3 \vspace{-3pt}\\
%     \hspace{15pt}{\footnotesize \vdots}       
%   \end{array}
%   \right.
%   \)
%   }
%   \end{minipage}
% \end{center}
%}
}





% \frame{
%   \frametitle{\insertsection \quad \quad{\normalsize $\cdots$ \insertsubsection}}
% \uncover<1->{
% {\small $\gamma$ : $\tilde{M}^5_2$のnull曲線, \, $A:=\dot{\gamma}$.} \vspace{0pt} 

% \hspace{-225pt}
% \scalebox{0.8}{\input{pic_M^5_2Frenet-2_Mensetsu}}
% }

% \uncover<2->{
% \begin{table}
% \begin{tabular}{c|c}
% \begin{minipage}{0.42\textwidth}
% %\underline{{\footnotesize {\bf Case 1} or {\bf Case 2} }}
% \vspace{-3pt}
% \begin{block}{{\small Type (\rnum{1})\, (H. Kobayashi)}}
%   \vspace{-5pt}
% {\footnotesize
% \begin{eqnarray}\label{joken5}
% {\renewcommand\arraystretch{0.3}
% {\arraycolsep = 0.5mm
% \hspace{-30pt}
% \begin{array}{l}
% A \\ B\\ C\\ Z_1\\ Z_2
% \end{array}
% \,
% \left(
% \begin{array}{ccccc}
% 0 & -1&   &   &   \\
% -1& 0 &   &   &   \\
%   &   & \varepsilon_C &            & \\
%   &   &            & \varepsilon_1 & \\
%   &   &            &            & \varepsilon_2
% \end{array}
% \right)
% }}
% \end{eqnarray}
% \vspace{-15pt}
% \begin{eqnarray}\label{joken5-2}
% \left \{
% \begin{array}{l}
% \dot{A}=k_1C\\
% \dot{C}=k_2A+\varepsilon_C k_1B\\
% \dot{B}=\varepsilon_C k_2 C+k_3Z_1+\varepsilon \gamma \\
% \dot{Z}_1=\varepsilon_1 k_3A+k_4Z_2 \\
% \dot{Z}_2=\varepsilon_C k_4Z_1
% \end{array}
% \right.
% \end{eqnarray}
% }
% \vspace{-9pt}
% \end{block}
% \end{minipage}
% }
% \uncover<3->{
% &
% \hspace{3pt}
% \begin{minipage}{0.42\textwidth}
% %\underline{{\footnotesize {\bf Case 3} }}
% \vspace{-3pt}
% \begin{block}{{\small Type (\rnum{2})\, (H. Kobayashi)}}
%   \vspace{-5pt}
% {\footnotesize
% \begin{eqnarray}\label{joken6}
% {\renewcommand\arraystretch{0.3}
% {\arraycolsep = 0.5mm
% \hspace{-30pt}
% \begin{array}{l}
% A \\ B\\ C\\ Z_1\\ Z_2
% \end{array}
% \left(
% \begin{array}{ccccc}
% 0 & -1&   &   &   \\
% -1& 0 &   &   &   \\
%   &   & 1 &   &   \\
%   &   &   & 0 & -1\\
%   &   &   & -1& 0 
% \end{array}
% \right)
% }}
% \end{eqnarray}
% \vspace{-15pt}
% \begin{eqnarray}\label{joken6-2}
% \left \{
% \begin{array}{l}
% \dot{A}=k_1C\\
% \dot{C}=k_2A+k_1B\\
% \dot{B}=k_2 C+Z_1+\varepsilon \gamma \\
% \dot{Z}_1=hZ_1 \\
% \dot{Z}_2=-A-hZ_2
% \end{array}
% \right.
% \end{eqnarray}
% }
% \vspace{-9pt}
% \end{block}
% \end{minipage}
% }
% \end{tabular}
% \end{table}
% }


\subsection{in \(\mathbb{S}^5_2\) or \(\mathbb{H}^5_2\)}

\frame{
  \frametitle{\insertsection \quad{\normalsize $\cdots$ \insertsubsection}}
%   \begin{itemize}
%   \item $s=1$のとき \quad $\gamma$ : $\tilde{M}^n_1$のnull曲線, \, $A:=\dot{\gamma}$. 
%   \end{itemize}
% \medskip
% \hspace{-3pt}
% \scalebox{0.88}{\input{pic_M^n_1Frenet}}
% \vspace{28pt}\\
% \pause

null曲線 \(\gamma \subset \tilde{M}^m_2\)に沿うFrenet型フレーム場, \, 
\vspace{7pt}

$A:=\dot{\gamma}$ 
\medskip
\pause
\hspace{-265pt}
\scalebox{0.9}{\input{pic_M^n_2Frenet-4_Mensetsu}}
\vspace{5pt}
\pause

\uncover<3->{\hspace{200pt}{\small 別証明}}
\vspace{-8pt}

\hspace{30pt}
\(
    \left\{
  \begin{aligned}
    & \text{\, Type (\rnum {1}) \hspace{-0pt} : Cartan frame field \uncover<3->{\hspace{37pt} {\small \(\cdots\) H. Kobayashi}}}\\
    & \text{\, Type (\rnum {2}) : bi-null Cartan frame field \uncover<3->{\hspace{7pt} {\small \(\cdots\) M. Sakaki--A. U\c{c}um--K. \.{I}larslan}}}\\
    & \text{\, Type (\rnum {3}) \uncover<3->{\hspace{128pt} {\small \(\cdots \,m=5\)のとき，H. Kobayashi}}}
  \end{aligned}
    \right.
  \)
}











\subsection{in \(\mathbb{S}^5_2\) or \(\mathbb{H}^5_2\)}

\frame{
  \frametitle{\insertsection \quad \quad{\normalsize $\cdots$ \insertsubsection}}
  \vspace{-3pt}

\begin{alertblock}{Main Result 4 (H. Kobayashi)}
      {\it 
      % \hspace{30pt}
      % \(\gamma\) \hspace{29pt}: \textcolor{red}{\(\mathbb{S}^5_2\) or \(\mathbb{H}^5_2\)}のnull曲線
      % \vspace{3pt}

      % \((A, B, C, Z_1, Z_2)\) : $\gamma$に沿う\( \mathbb{S}^5_2\) or \( \mathbb{H}^5_2\)のCartan frame field \, s.t.\, \(k_2\) : const
      % \vspace{10pt}
      % \pause

      % \(M\)が次の条件を満たすとする：\vspace{-5pt}
      % \[
      %   (*)\quad
      %   \left\{
      % \begin{aligned}
      %   & M \text{\,: \(\gamma\)と\(\gamma\)に沿うFrenet型フレーム場から構成されるnull scroll}\\
      %   & M \text{\,: \(2\)次元非退化ローレンツ線織面}
      % \end{aligned}
      %   \right.
      %   \vspace{-10pt}
      % \]
      % % \[
      % % \text{s.t.} \,
      % % \left\{  
      % % \begin{array}{l}
      % % B \text{\,: nullベクトル場,}\\
      % % \langle A, B \rangle =-1 ,\,\, \langle B, C \rangle =0.
      % % \end{array}
      % % \right.
      % \pause
      % \begin{center}
      % \(\Rightarrow\)\, \(M\)は\, \({\bf x}(s, t)=\gamma(s)+tB(s)\)\, によりパラメータづけされる．
      % \end{center}
      % \vspace{-0pt}
      % % \[
      % %   \Rightarrow \quad
      % % M: \, {\bf x}(s, t):=\gamma(s)+tB(s)
      % %   \vspace{-0pt}
      % % \]
      % % \hspace{3pt} \(\tilde{H}\) \hspace{5pt}: \(M\)の平均曲率ベクトル場，
      % % \(H:=\tilde{H}/\Vert \tilde{H}\Vert\)
      % % \vspace{3pt}
      % \pause

    % \hspace{43pt}
    % \begin{table}
    % \begin{tabular}{ll}
    %   \begin{minipage}{0.45\textwidth}
    %{\small
    \hspace{6pt}\(M\)\hspace{5pt} : \, \textcolor{red}{\(\mathbb{S}^5_2\) or \(\mathbb{H}^5_2\)}におけるnull scroll \vspace{3pt}

    \hspace{6pt}\(H\)\hspace{5pt} : \, \(M\)の平均曲率ベクトル場\vspace{3pt}
    
    \hspace{4pt}\(A_H\)\hspace{2pt} : \, \(M\)の\(H\)方向の形作用素\vspace{3pt}

    \(P(x)\) : \, \(A_H\)の最小多項式\vspace{7pt}
   % }
  %     \end{minipage}
  %     &
  %     \begin{minipage}{0.5\textwidth}
  %       {\small 
  %       \begin{itembox}{Note}
  %     \begin{itemize}
  %     \item null曲線\(\gamma\)と\(\gamma\)に沿うFrenet型フレーム場から構成される線織面
  %     \item \(2\)次元非退化ローレンツ曲面
  %     \item 形作用素は対角化不可能，実固有値
  %     \end{itemize}
  %   \end{itembox}
  %       }
  %     \end{minipage}
  %   \end{tabular}
  % \end{table}

    \(M\)がgeneralized B-scrollのとき，\(P(x)\)は次のいずれかである：
    
    \vspace{-3pt}
    \begin{center}
      \begin{minipage}{0.7\textwidth}
    \begin{enumerate}
    \uncover<2->{\item \(Z_1\) : non-null  \hspace{5pt}$\Rightarrow$\hspace{5pt}  \( P(x)=(x-(\textcolor<4->{red}{\varepsilon_C k_2^2+\varepsilon_1 k_3^2}) )^2\)}  \\
    \uncover<3->{\item \(Z_1\) : null  \hspace{26pt}$\Rightarrow$ \hspace{2pt} \( P(x)=(x-\textcolor<4->{red}{k_2^2})^2\)}
    \end{enumerate}
  \end{minipage}
    \end{center}
      %\vspace{-2pt}
  }
  \vspace{-5pt}
\end{alertblock}
\vspace{-6pt}
\begin{table}
\begin{tabular}{cc}
\begin{minipage}{0.42\textwidth}
%\underline{{\footnotesize {\bf Case 1} or {\bf Case 2} }}
\vspace{-3pt}
\uncover<2->{
\begin{block}{{\small Type (\rnum{1})\, \(\cdots Z_1\)がnon-nullのとき}}
  \vspace{-8pt}
{\footnotesize
% \begin{eqnarray}\label{joken5}
% {\renewcommand\arraystretch{0.3}
% {\arraycolsep = 0.5mm
% \hspace{-30pt}
% \begin{array}{l}
% A \\ B\\ C\\ Z_1\\ Z_2
% \end{array}
% \,
% \left(
% \begin{array}{ccccc}
% 0 & -1&   &   &   \\
% -1& 0 &   &   &   \\
%   &   & \varepsilon_C &            & \\
%   &   &            & \varepsilon_1 & \\
%   &   &            &            & \varepsilon_2
% \end{array}
% \right)
% }}
% \end{eqnarray}
%\vspace{-15pt}
\begin{eqnarray}\label{joken5-2}
\left \{
\begin{array}{l}
\dot{A}=k_1C\\
\dot{C}=\textcolor<4->{red}{k_2}A+\varepsilon_C k_1B\\
\dot{B}=\textcolor<4->{red}{\varepsilon_C k_2} C+\textcolor<4->{red}{k_3}Z_1+\varepsilon \gamma \\
\dot{Z}_1=\textcolor<4->{red}{\varepsilon_1 k_3}A+k_4Z_2 \\
\dot{Z}_2=\varepsilon_C k_4Z_1
\end{array}
\right.
\end{eqnarray}
}
\vspace{-9pt}
\end{block}
}
\end{minipage}
&
\hspace{3pt}
\begin{minipage}{0.42\textwidth}
%\underline{{\footnotesize {\bf Case 3} }}
\vspace{-3pt}
\uncover<3->{
\begin{block}{{\small Type (\rnum{3})\, \(\cdots Z_1\)がnullのとき}}
  \vspace{-8pt}
{\footnotesize
% \begin{eqnarray}\label{joken6}
% {\renewcommand\arraystretch{0.3}
% {\arraycolsep = 0.5mm
% \hspace{-30pt}
% \begin{array}{l}
% A \\ B\\ C\\ Z_1\\ Z_2
% \end{array}
% \left(
% \begin{array}{ccccc}
% 0 & -1&   &   &   \\
% -1& 0 &   &   &   \\
%   &   & 1 &   &   \\
%   &   &   & 0 & -1\\
%   &   &   & -1& 0 
% \end{array}
% \right)
% }}
% \end{eqnarray}
%\vspace{-15pt}
\begin{eqnarray}\label{joken6-2}
\left \{
\begin{array}{l}
\dot{A}=k_1C\\
\dot{C}=\textcolor<4->{red}{k_2}A+k_1B\\
\dot{B}=\textcolor<4->{red}{k_2} C+Z_1+\varepsilon \gamma \\
\dot{Z}_1=hZ_1 \\
\dot{Z}_2=-A-hZ_2
\end{array}
\right.
\end{eqnarray}
}
\vspace{-9pt}
\end{block}
}
\end{minipage}
\end{tabular}
\end{table}
}




% \section{今後の課題}


% \frame{
%   \frametitle{\insertsection}
% \underline{\quad$\tilde{M}^n_2$の形作用素\quad} \\
% \vspace{10pt}
% {\footnotesize
% \textcolor{red}{
% $
% \left(
% {\renewcommand\arraystretch{1}
% {\arraycolsep = 1mm
% \begin{array}{cc|ccc}
% \lambda & 0 & & & \\
%  1 & \lambda &  & \\ \hline
%  & &  \lambda &   & \\
%  & &  & \ddots & \\
%  & &    &   & \lambda 
% \end{array}
% }}
% \right)
% $
% }
% }
%  ,\quad
% {\footnotesize
% $
% \left(
% {\renewcommand\arraystretch{0.5}
% {\arraycolsep = 1mm
% \begin{array}{cc|cc|ccc}
% \lambda & 0 & & & & & \\
% 1 & \lambda & & & & & \\ \hline
%  & & \lambda & 0 & & & \\
%  & & 1 & \lambda & & & \\ \hline
%  & &    &   & \lambda & & \\
%  & &    &   &   & \ddots & \\
%  & &    &   &   &        & \lambda
% \end{array}
% }}
% \right)
% $
% }
%  ,\quad
% {\footnotesize
% $
% \left(
% {\renewcommand\arraystretch{0.5}
% {\arraycolsep = 1mm
% \begin{array}{ccc|ccc}
% \lambda &  &  &  & & \\
%  1&  \lambda &  &  & & \\
%  & 1& \lambda  &  & & \\ \hline
%  &  &  & \lambda &  & \\
%  &  &  &   & \ddots &  \\
%  &  &  &   &   & \lambda  
% \end{array}
% }}
% \right)
% $
% }
%  ,\quad
% \vspace{10pt} \\
% {\footnotesize
% $
% \left(
% {\renewcommand\arraystretch{0.8}
% {\arraycolsep = 1mm
% \begin{array}{ccc|ccc|ccc}
% \lambda &         &         &   &   &   &   &        &   \\
% 1       & \lambda &         &   &   &   &   &        &   \\
%         &  1      & \lambda &   &   &   &   &        &   \\ \hline
%         &         &         & \lambda &         &  &   &        &   \\ 
%         &         &         &   1     & \lambda &   &   &        &   \\
%         &         &         &         &   1     & \lambda &   &        &   \\ \hline
%  &   &   &   &   &   & \lambda &        &   \\
%  &   &   &   &   &   &   & \ddots &   \\
%  &   &   &   &   &   &   &        & \lambda \\
% \end{array}
% }}
% \right)
% $ ,\quad
% {\footnotesize
% $
% \left(
% {\renewcommand\arraystretch{0.5}
% {\arraycolsep = 1mm
% \begin{array}{cccc|ccc}
% \lambda & 0 & & & & & \\
% 1 & \lambda & & & & & \\
%  & 1 & \lambda & 0 & & & \\
%  & & 1 & \lambda & & & \\ \hline
%  & &    &   & \lambda & & \\
%  & &    &   &   & \ddots & \\
%  & &    &   &   &        & \lambda
% \end{array}
% }}
% \right)
% $
% }
%  ,\quad
% $
% \left(
% {\renewcommand\arraystretch{0.5}
% {\arraycolsep = 1mm
% \begin{array}{ccccc|ccc}
% \lambda&          &           &          &   &   &   &          \\
% 1      &  \lambda &           &          &   &          &          &   \\
%        &      1   & \lambda   &          &   &   &   &          \\
%        &          &  1        &  \lambda &   &  &   &          \\
%        &          &           &    1     & \lambda &   &          &   \\ \hline
%        &          &           &          &      & \lambda &        &   \\
%        &          &           &          &   &     & \ddots &   \\
%        &          &           &          &   &   &   &         \lambda  
% \end{array}
% }}
% \right)
% $
% }
% }


\end{document}